\documentclass[12pt,a4paper,onecolumn]{article}
\title{Digest of The New York Times}
\author{The New York Times}
% \definecolor{dark-green}{rgb}{0,0.25,0}
% \usepackage[utf8]{inputenc}

\usepackage{fontspec,xunicode,xltxtra}
\setmainfont[Ligatures=Required,Mapping=tex-text]{Georgia}
% \setmainfont[Ligatures=Required,Mapping=tex-text]{Palatino Linotype}
\setsansfont[Ligatures=Required,Mapping=tex-text]{Myriad Pro}
\setmonofont[Ligatures=Required,Mapping=tex-text]{Courier New}
\usepackage{xeCJK}
\setCJKmainfont{宋体}
\setCJKmonofont{黑体}

% \usepackage{indentfirst}


\usepackage[dvips,dvipsnames,svgnames]{xcolor}
\definecolor{light-gray}{gray}{0.95}

%生成PDF的链接
\usepackage[dvipdfmx]{hyperref}
\hypersetup{
    bookmarks=true,         % show bookmarks bar?
    bookmarksopen=true,
    pdfpagemode=UseNone,    % options: UseNode, UseThumbs, UseOutlines, FullScreen
    pdfstartview=FitH,
    pdfborder=1,
    pdfhighlight=/P,
    pdfauthor={wuxch},
    unicode=true,           % non-Latin characters in Acrobat’s bookmarks
    colorlinks,             % false: boxed links; true: colored links
    linkcolor=blue,         % color of internal links
    citecolor=blue,        % color of links to bibliography
    filecolor=magenta,      % color of file links
    urlcolor=cyan           % color of external links
}
% 使用目录
\usepackage{makeidx}
% 使用图形
\usepackage{graphicx}


\usepackage{titletoc}
\dottedcontents{section}[1.5em]{}{2.3em}{1pc}

% % 页眉页脚的设定
% \usepackage{fancyhdr}
% \pagestyle{fancy}
% \lhead{\includegraphics[height=0.2in]{zte_logo}}

% \usepackage{footmisc}
\usepackage{pfnote}
\usepackage{fnpos}

% % 嵌入的代码显示
% \usepackage{listings}
% \lstset{language=C++, breaklines, extendedchars=false}
% \lstset{keywordstyle=\color{blue},
%         commentstyle=\color{SeaGreen},
%         stringstyle=\ttfamily,
%         showstringspaces=false,
%         backgroundcolor=\color{light-gray}}
% \lstset{basicstyle=\small}

\usepackage{booktabs}

\usepackage{tipa}
\usepackage{phonetic}
\usepackage[gen]{eurosym}

% \usepackage{CJK}
\usepackage{makeidx}
% \usepackage{glossary}

\usepackage{fancyhdr}

\usepackage{lastpage}
\cfoot{\textsf{\thepage\ of \pageref*{LastPage}}}

\pagestyle{fancy}
\fancyhead{}
\renewcommand{\headrulewidth}{0pt}
\pagenumbering{arabic}

\usepackage{graphicx}
\usepackage{picinpar}

\setlength{\parindent}{0in}
\setlength{\parskip}{2.5ex plus4mm minus3mm}
% \setlength{\parskip}{.3\baselineskip}

\setcounter{section}{0}

% bibtex
\usepackage[super]{natbib}
% \usepackage[globalcitecopy,labelstoglobalaux]{bibunits}
% \makeatletter
% \renewcommand{\@bibunitname}{\jobname.\the\@bibunitauxcnt}
% \makeatother

% \renewcommand{\refname}{缩略语}
% \renewcommand{\refname}{}

\makeatletter                   % Reference list option change
\renewcommand\@biblabel[1]{#1:  } % from [1] to 1: 
\makeatother                    %

\newcommand{\citenumfont}[1]{\textsf{#1}}

% \usepackage{tipa}
% \usepackage{phonetic}

\newcommand{\phonetic}[1]{{\fontspec[Colour=CC00CC]{Microsoft Sans Serif} #1}}


%%% Local Variables:
%%% TeX-master: "medical_trivia"
%%% End:

\makeindex
\begin{document}
\date{}

% \maketitle
\tableofcontents
\pagebreak
\section{In China, G.M. Remains a Driving Force}

SHANGHAI -- Ford may be standing taller than General Motors in Detroit these days -- flush with cash
while its rival is forced to go repeatedly to Washington, hat in hand, seeking government bailouts.
But in China the tables are turned.

G.M. is a powerful presence here with 8 to 10 percent of the market for cars, minivans and sport
utility vehicles, making it the second-largest automaker in China for such vehicles, passed only by
Volkswagen. One of G.M.'s local joint ventures, Wuling, dominates the sale of bare-bones pickups and
vans, hugely popular in rural areas, with nearly half the market.

And in discussions about a possible bankruptcy to break G.M. into a ``good G.M.'' of viable
businesses and a ``bad G.M.'' composed of the rest of the automaker, the company's Chinese
operations are consistently mentioned as a candidate for the good G.M.

Unlike the gas guzzlers churned out by G.M.'s North American operation, the company's China division
has emphasized fuel-sipping models in the last decade. Two weeks ago, G.M. set a target of roughly
doubling sales in China within five years, to two million vehicles a year.

By contrast, Ford languishes in 10th place with 3 percent of the market.

G.M.'s recent fortunes received a lift from the Chinese government's decision last autumn to cut
taxes to 1 percent on cars with engines of 1.6 liters or less while raising them to as much as 40
percent for cars, minivans and sport utility vehicles with larger engines.

And after the government introduced a subsidy of \$730 last month for vehicle buyers in rural areas,
G.M.'s Wuling operation can barely build minivans fast enough.

Only one of the four models Ford sells in China, the Fiesta, has engines small enough to qualify for
the recent tax break. Fiesta sales have been strong while sales of other models have been weak.

Ford's joint venture partner for making sedans in China, Changan Automobile Group, is the main rival
for G.M.'s joint venture in small minivans -- but Ford set up its relationship with Changan many
years ago to exclude such light commercial vehicles, so Changan makes them on its own.

Nick Reilly, the president of G.M.'s Asian and Pacific region, said at the Shanghai auto show this
week that G.M.'s Chinese operations generated enough cash to cover their own investment costs.

Beyond the recent tax incentives and subsidies, G.M. is ahead in China for two other reasons. The
company bought the assets of bankrupt Daewoo of South Korea in 2002 and turned them into its
worldwide center for small-car design, with an emphasis on supplying the Asian market.

Ford has continued to rely on Europe to design its small cars, although John Parker, Ford's
executive vice president for Asia, the Pacific and Africa, said Ford would look for ways to do more
engineering in China as well.

G.M.'s other advantage is that its Chinese operations were highly independent of the rest of the
corporation for many years. Philip F.~Murtaugh, the executive who built G.M. China, emphasized fuel
economy and affordability as the company's North American operations were betting heavily on
S.U.V.'s. In fact, G.M. became so big and successful in China that the company's management stripped
the subsidiary of its independence in early 2005 and put vehicle design, parts purchasing and other
functions under more supervision from Detroit. (Mr.~Murtaugh resigned in protest; he has since
worked stints at the Shanghai Automotive Industry Corporation and Chrysler and has been unemployed
and looking for work since December.)

Ford has taken what it has sometimes described as a ``measured'' approach to the Chinese market --
only to be caught by surprise as the Chinese market surpassed the American market in sales during
the first quarter.

Mr.~Parker said that the manufacturer was reviewing its offerings in the Chinese market in light of
government policies aimed at pushing the market toward fuel-efficient models that are also
affordable.

But Ford has no intention to enter the market for light commercial vehicles in China, he added.
Costing as little as \$5,000, these vehicles share few parts with cars.

Ford sold 306,000 vehicles in China last year compared with 1.9 million from its core Ford, Lincoln
and Mercury brands in the United States. G.M. sold 1.09 million vehicles in China last year, a third
of its sales in the United States.

G.M.'s vehicles in China range from bare-bones pickups for less than \$5,000 to Cadillacs costing
more than \$100,000.

The weakest of the Detroit automakers here is Chrysler, which does not even rank among the top 20
automakers in China by sales . When Daimler sold most of Chrysler to Cerberus nearly two years ago,
a long argument ensued over the fate of a factory on the outskirts of Beijing that was making both
Mercedes-Benz and Chrysler cars.

Daimler ended up keeping the factory, leaving Chrysler dependent on imports, which are subject to a
25 percent tariff.

A small stage was erected as a backdrop for Chrysler's exhibit of existing models, with a row of
fake cactus and a toy wind turbine slowly rotating at the top of a 15-foot-high tower.

But Chrysler has no vehicles on the market in China with engines of 1.6 liters or less; the Chrysler
vehicle with the smallest engine on display at the show was a Dodge Caliber compact with a 2-liter
engine. In front of the fake cactus was a Jeep Commander with a 5.7-liter engine.

Michael Dunne, an auto analyst in Shanghai for J.D. Power \& Associates, said that Chinese car
buyers were aware of the difficulties of the Detroit automakers in the United States because of
heavy media coverage of them in China, but tended to see their operations in China as distinct.
Executives at Nissan and at three Chinese automakers -- Geely, Great Wall Motor and New-Power
Automobile Technology -- all said that they saw little effect from the American difficulties on the
sales here of G.M. and Ford.

Indeed, Detroit's difficulties may have a side benefit of shielding the American market from Chinese
competition for a little while longer. Companies like Great Wall are wary of starting to export to
the United States at such a politically and financially worrisome time.

``I don't think this year or the next year is a good time to launch a new brand, because the Big
Three have headaches and the U.S.~government is nervous,'' said Steven Wang, Great Wall's marketing
director.

\section{Basketball Prospect Leaving High School to Play in Europe}
SAN DIEGO -- Jeremy Tyler, a 6-foot-11 high school junior whom some consider the best American big
man since Greg Oden, says he will be taking a new path to the N.B.A. He has left San Diego High
School and said this week that he would skip his senior year to play professionally in Europe.

Tyler, 17, would become the first United States-born player to leave high school early to play
professionally overseas. He is expected to return in two years, when he is projected to be a top
pick, if not the No.~1 pick, in the 2011 N.B.A. draft.

Tyler, who had orally committed to play for Rick Pitino at Louisville, has yet to sign with an agent
or a professional team. His likely destination is Spain, though teams from other European leagues
have shown interest. A spokesman for Louisville said the university could not comment about Tyler.

``Nowadays people look to college for more off-the-court stuff versus being in the gym and getting
better,'' Tyler said. ``If you're really focused on getting better, you go play pro somewhere. Pro
guys will get you way better than playing against college guys.''

His decision is perhaps the most important one since Kevin Garnett jumped straight to the N.B.A.
from high school in 1995. Garnett was the No.~5 pick in the N.B.A. draft and ushered in a generation
of preps-to-pros stars like Kobe Bryant, LeBron James and Dwight Howard. A minimum-age rule passed
for the 2006 draft cut off that route, essentially forcing players to spend at least one year in
college.

But Brandon Jennings, a point guard from Los Angeles, became the first player to graduate from high
school, skip college and play professionally in Europe. (Whether Jennings would have qualified
academically to play at Arizona, where he had signed a letter of intent, is unknown.) He is in his
first season with Lottomatica Virtus Roma in Italy and is projected as a high pick in the N.B.A.
draft in June.

Tyler took Jennings's path and added a compelling twist, perhaps opening the door for other elite
high school basketball players to follow.

Sonny Vaccaro, a former sneaker company executive, orchestrated Jennings's move and has guided Tyler
and his family through the process.

``It's significant because it shows the curiosity for the American player just refusing to accept
what he's told he has to do,'' Vaccaro said. ``We're getting closer to the European reality of a
professional at a young age. Basically, Jeremy Tyler is saying, 'Why do I have to go to high
school?' ''

Vaccaro said he was unsure how much money Tyler would make, though it will most likely be less than
the \$1.2 million Jennings made in a combination of salary and endorsements this season. Vaccaro
said Tyler would make a six-figure salary, noting that the economic crisis in Europe could hurt his
earnings.

Vaccaro made his name by signing Michael Jordan to Nike in the mid-1980s and has advised numerous
elite players over the years. ``I believe he'll be a 10-time All-Star with his ability,'' Vaccaro
said of Tyler.

For now, Tyler said he was working out eight to nine hours a day at local gyms and schools and
spending two and a half hours a day to get his high school diploma online. He said he had his ``ups
and downs'' in the classroom but was on track to qualify academically to play in college.

Tyler said that his game stagnated by playing high school basketball here and that he was frustrated
by the rules of the California Interscholastic Federation and the National Collegiate Athletic
Association.

He averaged 28.7 points during a tumultuous junior season in which his team went 15-11. Two of his
coaches were fired in part because of their roles in trying to bring in three star players to
complement Tyler. Kenny Roy, the former head coach, denied any wrongdoing.

But the three transfers were ruled ineligible, leaving Tyler the centerpiece of a team whose next
biggest player was 6-2.

``It was boring and I wasn't getting better,'' Tyler said. ``Each game was the same thing. I was
getting triple-teamed and getting hacked. After each game I'd have scratches and bruises up and down
my arms from getting triple-teamed. It just wasn't for me.''

Spain is the most likely option for him because of the structure of its basketball league. Vaccaro
said Tyler could shuttle between his professional team and its development squad for 18-year-olds.
So if Tyler, who turns 18 in June, struggles early against professional men, he will be able to play
and practice with the younger team.

Tyler, who was raised primarily by his father, James, and his uncle, Maurice, will move to Europe
with his older brother, James. A rotation of relatives plan to visit. Tyler's father said he had
reservations at first about sending his son overseas, knowing there would be a strong reaction in
the community.

``You know what? It's just a job,'' the elder James Tyler said. ``He gets a chance to work the job
of his dreams and have fun doing it. There's kids leaving to go to Iraq at age 18. They took a job
to serve their country. He's going to play ball.''

Along the way, Tyler may just become a trailblazer for other high school basketball prodigies.

``This is what I want to do,'' he said. ``It will help me get better and grow up.''

\section{Japan Pays Foreign Workers to Go Home}

HAMAMATSU, Japan -- Rita Yamaoka, a mother of three who immigrated from Brazil, recently lost her
factory job here. Now, Japan has made her an offer she may not be able to refuse.

The government will pay thousands of dollars to fly Mrs.~Yamaoka; her husband, who is a Brazilian
citizen of Japanese descent; and their family back to Brazil. But in exchange, Mrs.~Yamaoka and her
husband must agree never to seek to work in Japan again.

``I feel immense stress. I've been crying very often,'' Mrs.~Yamaoka, 38, said after a meeting where
local officials detailed the offer in this industrial town in central Japan.

``I tell my husband that we should take the money and go back,'' she said, her eyes teary. ``We
can't afford to stay here much longer.''

Japan's offer, extended to hundreds of thousands of blue-collar Latin American immigrants, is part
of a new drive to encourage them to leave this recession-racked country. So far, at least 100
workers and their families have agreed to leave, Japanese officials said.

But critics denounce the program as shortsighted, inhumane and a threat to what little progress
Japan has made in opening its economy to foreign workers.

``It's a disgrace. It's cold-hearted,'' said Hidenori Sakanaka, director of the Japan Immigration
Policy Institute, an independent research organization.

``And Japan is kicking itself in the foot,'' he added. ``We might be in a recession now, but it's
clear it doesn't have a future without workers from overseas.''

The program is limited to the country's Latin American guest workers, whose Japanese parents and
grandparents emigrated to Brazil and neighboring countries a century ago to work on coffee
plantations.

In 1990, Japan -- facing a growing industrial labor shortage -- started issuing thousands of special
work visas to descendants of these emigrants. An estimated 366,000 Brazilians and Peruvians now live
in Japan.

The guest workers quickly became the largest group of foreign blue-collar workers in an otherwise
immigration-averse country, filling the so-called three-K jobs (kitsui, kitanai, kiken -- hard,
dirty and dangerous).

But the nation's manufacturing sector has slumped as demand for Japanese goods evaporated, pushing
unemployment to a three-year high of 4.4 percent. Japan's exports plunged 45.6 percent in March from
a year earlier, and industrial production is at its lowest level in 25 years.

New data from the Japanese trade ministry suggested manufacturing output could rise in March and
April, as manufacturers start to ease production cuts. But the numbers could have more to do with
inventories falling so low that they need to be replenished than with any increase in demand.

While Japan waits for that to happen, it has been keen to help foreign workers leave, which could
ease pressure on domestic labor markets and the unemployment rolls.

``There won't be good employment opportunities for a while, so that's why we're suggesting that the
Nikkei Brazilians go home,'' said Jiro Kawasaki, a former health minister and senior lawmaker of the
ruling Liberal Democratic Party.

``Nikkei'' visas are special visas granted because of Japanese ancestry or association.

Mr.~Kawasaki led the ruling party task force that devised the repatriation plan, part of a wider
emergency strategy to combat rising unemployment.

Under the emergency program, introduced this month, the country's Brazilian and other Latin American
guest workers are offered \$3,000 toward air fare, plus \$2,000 for each dependent -- attractive
lump sums for many immigrants here. Workers who leave have been told they can pocket any amount left
over.

But those who travel home on Japan's dime will not be allowed to reapply for a work visa. Stripped
of that status, most would find it all but impossible to return. They could come back on three-month
tourist visas. Or, if they became doctors or bankers or held certain other positions, and had a
company sponsor, they could apply for professional visas.

Spain, with a unemployment rate of 15.5 percent, has adopted a similar program, but immigrants are
allowed to reclaim their residency and work visas after three years.

Japan is under pressure to allow returns. Officials have said they will consider such a
modification, but have not committed to it.

``Naturally, we don't want those same people back in Japan after a couple of months,'' Mr.~Kawasaki
said. ``Japanese taxpayers would ask, 'What kind of ridiculous policy is this?' ''

The plan came as a shock to many, especially after the government introduced a number of measures in
recent months to help jobless foreigners, including free Japanese-language courses, vocational
training and job counseling. Guest workers are eligible for limited cash unemployment benefits,
provided they have paid monthly premiums.

``It's baffling,'' said Angelo Ishi, an associate professor in sociology at Musashi University in
Tokyo. ``The Japanese government has previously made it clear that they welcome Japanese-Brazilians,
but this is an insult to the community.''

It could also hurt Japan in the long run. The aging country faces an impending labor shortage. The
population has been falling since 2005, and its working-age population could fall by a third by
2050. Though manufacturers have been laying off workers, sectors like farming and care for the
elderly still face shortages.

But Mr.~Kawasaki said the economic slump was a good opportunity to overhaul Japan's immigration
policy as a whole.

``We should stop letting unskilled laborers into Japan. We should make sure that even the three-K
jobs are paid well, and that they are filled by Japanese,'' he said. ``I do not think that Japan
should ever become a multi-ethnic society.''

He said the United States had been ``a failure on the immigration front,'' and cited extreme income
inequalities between rich Americans and poor immigrants.

At the packed town hall meeting in Hamamatsu, immigrants voiced disbelief that they would be barred
from returning. Angry members of the audience converged on officials. Others walked out of the
meeting room.

``Are you saying even our children will not be able to come back?'' one man shouted.

``That is correct, they will not be able to come back,'' a local labor official, Masahiro Watai,
answered calmly.

Claudio Nishimori, 30, said he was considering returning to Brazil because his shifts at a
electronics parts factory were recently reduced. But he felt anxious about going back to a country
he had left so long ago.

``I've lived in Japan for 13 years. I'm not sure what job I can find when I return to Brazil,'' he
said. But his wife has been unemployed since being laid off last year and he can no longer afford to
support his family.

Mrs.~Yamaoka and her husband, Sergio, who settled here three years ago at the height of the export
boom, are undecided. But they have both lost jobs at auto factories. Others have made up their minds
to leave. About 1,000 of Hamamatsu's Brazilian inhabitants left the city before the aid was even
announced. The city's Brazilian elementary school closed last month.

``They put up with us as long as they needed the labor,'' said Wellington Shibuya, who came six
years ago and lost his job at a stove factory in October. ``But now that the economy is bad, they
throw us a bit of cash and say goodbye.''

He recently applied for the government repatriation aid and is set to leave in June.

``We worked hard; we tried to fit in. Yet they're so quick to kick us out,'' he said. ``I'm happy to
leave a country like this.''

\section{IPhone Sales Push Up Apple's Profits}

Apple's second-quarter profit, driven by strong sales of iPhones and iPods, beat Wall Street's
expectations despite the gloomy economy.

The company, based in Cupertino, Calif., said on Wednesday that continued demand for its portable
devices kept sales from dipping substantially during the typically slow months after the holiday
shopping season.

Despite a 3 percent slump in sales of the company's Mac computers, Peter Oppenheimer, Apple's chief
financial officer, called the quarter ``the best nonholiday earnings in the company's
history.''Apple reported that its net profit jumped 15 percent, to \$1.21 billion, or \$1.33 a share,
from \$1.05 billion, or \$1.16 a share, in the same quarter a year ago.

Quarterly revenue rose to \$8.16 billion, from \$7.51 billion last year. Analysts had expected Apple
to announce revenue of \$7.9 billion and a profit of \$1.09 a share, according to a survey conducted
by Thomson Reuters.

Analysts, including A.~M. Sacconaghi of Sanford C.~Bernstein \& Company, said the most impressive
figure was the company's 36.4 percent gross margin, which Apple said was helped by the falling cost
of some components and strong software sales.

Mr.~Sacconaghi also remarked on the high spirits of Mr.~Oppenheimer and the company's chief
operating officer, Timothy D.~Cook, during a conference call with analysts and shareholders.

``Their enthusiasm about their product portfolio is something investors care deeply about,''
Mr.~Sacconaghi said. ``Even though they weren't specific, that level of excitement potentially
foreshadows exciting products to come. It is something investors care deeply about and get excited
over.''

Apple watchers anticipate that the company will unveil new iPhones at its Worldwide Developers
Conference in June and perhaps a tablet device later in the year. Discussing rumors that the company
would develop a so-called netbook -- a small, low-cost laptop with limited features -- Mr.~Cook
expressed little interest.

``When I look at what is being sold in the netbook space today, I see cramped keyboards, terrible
software, junky hardware, very small screens,'' he said. ``And just not a consumer experience, and
not something that we would put the Mac brand on, quite frankly.''

Mr.~Cook also said that the iPod Touch or iPhone were worthy alternatives to ``small computers that
do browsing and e-mail.''

When pressed about the condition of Steven P.~Jobs, Apple's chief executive, who is on a leave of
absence to deal with health matters, Mr.~Oppenheimer was upbeat. ``We look forward to Steve
returning to Apple at the end of June,'' he said.

The earnings news sent Apple's stock up 3.5 percent in after-hours trading. The earnings report
offered more proof that Apple's strategy of selling must-have digital accoutrements at a premium
price to people who can afford them is a smart, if counterintuitive, strategy for navigating the
recession.

``For a consumer company to be growing 9 percent in this economic environment is a testament that
the products are strong and the company is in touch with its customers,'' Mr.~Sacconaghi said.

Apple said it sold 2.22 million Macs in the quarter ended March 28, down from 2.3 million last year.
But it sold 3.79 million iPhones, up from 1.7 million a year ago, in part because of expanding sales
of the device around the globe.

Mr.~Cook said that the App Store for the iPhone, which offers software from third-party developers,
was within hours of approaching its one billionth download.

The planned summer update to the iPhone's operating system, which will include features like the
ability to make payments within applications and multimedia messaging, ``unleashes a whole new level
of innovation that keeps Apple years ahead of everyone else,'' he said.

As usual, the company was cautious in the guidance it offered Wall Street. The company projected
that it would have revenue of \$7.7 billion to \$7.9 billion and earnings of about \$1 a share in
the current quarter.

Analysts applauded the company's earnings report but cautioned against taking Apple's gains as a
sign of the wider industry's improving.

``Apple is a big company, but they are really addressing only the upper ends of their target
markets, and those are arguably the parts of the market that are most resistant to the economic
downturn,'' said Shaw Wu, a senior research analyst at Kaufman Brothers.

\section{Court Debates Strip Search of Student}

WASHINGTON — The United States Supreme Court spent an hour on Tuesday debating what middle school
students are apt to put in their underwear and what should be done about it.

Justice Stephen G.~Breyer, for instance, said it struck him as ``a logical thing'' that adolescents
seeking to hide pills ``will stick them in their underwear.''

Adam B.~Wolf, a lawyer with the American Civil Liberties Union, disagreed, invoking what he called
the ``ick factor.''

His client, Savana Redding, had been subjected to a strip search in 2003 by school officials in
Safford, Ariz. She was 13 and in eighth grade at the time.

The officials were acting on a tip from another student and were looking for prescription-strength
ibuprofen, a painkiller. They made Ms.~Redding strip to her underwear, shake her bra and pull aside
her panties. The officials, both female, found no pills.

``What this school official did,'' Mr.~Wolf said, referring to the male assistant principal who
ordered the search, ``was act on nothing more than a hunch -- if that -- that Savana was currently
concealing ibuprofen pills underneath her underpants for others' oral consumption.''

``I mean, there's a certain ick factor to this,'' Mr.~Wolf said.

Without intimating a view on the ickiness of what Mr.~Wolf had described, Chief Justice John
G.~Roberts Jr.~suggested that the law might treat different undergarments differently. ``The issue
here covers the brassiere as well,'' he said, ``which doesn't seem as outlandish as the
underpants.''

Justice Breyer elaborated on what children put in their underwear. ``In my experience when I was 8
or 10 or 12 years old, you know, we did take our clothes off once a day,'' he said. ``We changed for
gym, O.K.? And in my experience, too, people did sometimes stick things in my underwear.''

The courtroom rocked with laughter, and the justice grew a little flustered at having apparently
misspoken.

While Supreme Court arguments can often be bone-dry exercises in statutory exegesis and doctrinal
refinement, Tuesday's session was grounded in vivid facts: school snitches, drugs, underwear and
body cavities.

None of the lawyers had a particularly easy time of it. Matthew W.~Wright, representing the school
district, said that intimate searches should be allowed even for the most common over-the-counter
drugs.

``At some point it gets silly,'' Justice David H.~Souter said. ``Having an aspirin tablet does not
present a health and safety risk.''

Mr.~Wright did draw the line at searches of students' body cavities, but only on the practical
ground that school officials are not trained to conduct such searches. Mr.~Wright said there was no
legal obstacle to such a search.

David O'Neil, an assistant to the solicitor general representing the federal government, tried to
steer a middle course.

The Fourth Amendment had been violated, he said, because school officials did not have a reasonable
suspicion that Ms.~Redding had secreted drugs in her undergarments. But Mr.~O'Neil added that
Ms.~Redding should not be allowed to sue the assistant principal who ordered the search, because the
law was unclear at the time.

Justice Antonin Scalia challenged him on the first point.

``You search in the student's pack, you search the student's outer garments, and you have a
reasonable suspicion that the student has drugs,'' he said. ``Don't you have, after conducting all
these other searches, a reasonable suspicion that she has drugs in her underpants?''

``You've searched everywhere else,'' Justice Scalia said. ``By God, the drugs must be in her
underpants.''

Mr.~O'Neil said a more focused suspicion was required. ``Certainly there is no practice anywhere,
that I'm aware of, of hiding ibuprofen in underwear,'' he said.

On the other hand, if there is good reason to think a student is keeping contraband in her
underwear, Mr.~O'Neil said, ``you could go directly to that location.''

That surprised the chief justice. ``Oh, surely not,'' he said. ``You are saying if you have
reasonable suspicion that it's in the underwear, you shouldn't even bother searching the pack or the
pockets? You should go straight to the underwear? That can't be right.''

Mr.~Wolf, Ms.~Redding's lawyer, injected another new term into the court's lexicon. He said a search
may be appropriate if the school has evidence that a student makes a habit of ``crotching'' drugs.

Justice Souter may have summarized the mood of the court near the end of the argument in the case,
Safford Unified School District v. Redding, No.~08-479. Several justices appeared troubled by the
search, but also seemed loath to second-guess school officials confronted with a variety of
dangerous substances.

``My thought process,'' Justice Souter said, ``is I would rather have the kid embarrassed by a strip
search, if we can't find anything short of that, than to have some other kids dead because the stuff
is distributed at lunchtime and things go awry.''

\section{In Adopting Harsh Tactics, No Inquiry Into Their Past Use}

WASHINGTON -- The program began with Central Intelligence Agency leaders in the grip of an alluring
idea: They could get tough in terrorist interrogations without risking legal trouble by adopting a
set of methods used on Americans during military training. How could that be torture?

In a series of high-level meetings in 2002, without a single dissent from cabinet members or
lawmakers, the United States for the first time officially embraced the brutal methods of
interrogation it had always condemned.

This extraordinary consensus was possible, an examination by The New York Times shows, largely
because no one involved -- not the top two C.I.A. officials who were pushing the program, not the
senior aides to President George W.~Bush, not the leaders of the Senate and House Intelligence
Committees -- investigated the gruesome origins of the techniques they were approving with little
debate.

According to several former top officials involved in the discussions seven years ago, they did not
know that the military training program, called SERE, for Survival, Evasion, Resistance and Escape,
had been created decades earlier to give American pilots and soldiers a sample of the torture
methods used by Communists in the Korean War, methods that had wrung false confessions from
Americans.

Even George J.~Tenet, the C.I.A. director who insisted that the agency had thoroughly researched its
proposal and pressed it on other officials, did not examine the history of the most shocking method,
the near-drowning technique known as waterboarding.

The top officials he briefed did not learn that waterboarding had been prosecuted by the United
States in war-crimes trials after World War II and was a well-documented favorite of despotic
governments since the Spanish Inquisition; one waterboard used under Pol Pot was even on display at
the genocide museum in Cambodia.

They did not know that some veteran trainers from the SERE program itself had warned in internal
memorandums that, morality aside, the methods were ineffective. Nor were most of the officials aware
that the former military psychologist who played a central role in persuading C.I.A. officials to
use the harsh methods had never conducted a real interrogation, or that the Justice Department
lawyer most responsible for declaring the methods legal had idiosyncratic ideas that even the Bush
Justice Department would later renounce.

The process was ``a perfect storm of ignorance and enthusiasm,'' a former C.I.A. official said.

Today, asked how it happened, Bush administration officials are finger-pointing. Some blame the
C.I.A., while some former agency officials blame the Justice Department or the White House.

Philip D.~Zelikow, who worked on interrogation issues as counselor to Secretary of State Condoleezza
Rice in 2005 and 2006, said the flawed decision-making badly served Mr.~Bush and the country.

``Competent staff work could have quickly canvassed relevant history, insights from the best law
enforcement and military interrogators, and lessons from the painful British and Israeli
experience,'' Mr.~Zelikow said. ``Especially in a time of great stress, walking into this minefield,
the president was entitled to get the most thoughtful and searching analysis our government could
muster.''

After years of recriminations about torture and American values, Bush administration officials say
it is easy to second-guess the decisions of 2002, when they feared that a new attack from Al Qaeda
could come any moment.

If they shunned interrogation methods some thought might work, and an undetected bomb or bioweapon
cost thousands of lives, where would the moral compass point today? It is a question that still
haunts some officials. Others say that if they had known the full history of the interrogation
methods or been able to anticipate how the issue would explode, they would have advised against
using them.

This account is based on interviews with more than two dozen current and former senior officials of
the C.I.A., White House, Justice Department and Congress. Nearly all, citing the possibility of
future investigations, shared their recollections of the internal discussions of a classified
program only on condition of anonymity.

Leaked to the news media months after they were first used, the C.I.A.'s interrogation methods would
darken the country's reputation, blur the moral distinction between terrorists and the Americans who
hunted them, bring broad condemnation from Western allies and become a ready-made defense for
governments accused of torture. The response has only intensified since Justice Department legal
memos released last week showed that two prisoners were waterboarded 266 times and that C.I.A.
interrogators were ordered to waterboard one of the captives despite their belief that he had no
more information to divulge.

But according to many Bush administration officials, including former Vice President Dick Cheney and
some intelligence officers who are critics of the coercive methods, the C.I.A. program would also
produce an invaluable trove of information on Al Qaeda, including leads on the whereabouts of
important operatives and on terror schemes discussed by Al Qaeda. Whether the same information could
have been acquired using the traditional, noncoercive methods that the Federal Bureau of
Investigation and the military have long used is impossible to say, and former Bush administration
officials say they did not have the luxury of time to develop a more patient approach, given that
they had intelligence warnings of further attacks.

Michael V.~Hayden, who served as C.I.A. director for the last two years of the Bush administration,
devoted part of his last press briefing in January to defending the C.I.A. program. ``It worked,''
Mr.~Hayden insisted.

``I have said to all who will listen that the agency did none of this out of enthusiasm,'' he said.
``It did it out of duty. It did it with the best legal advice it had.''

A Program Takes Shape

When Mr.~Bush assigned the C.I.A. with the task of questioning high-level Qaeda captives in late
2001, the agency had almost no experience interrogating the kind of hostile prisoners it soon
expected to hold.

It had dozens of psychiatrists, psychologists, polygraphists and operations officers who had
practiced the arts of eliciting information and assessing truthfulness. Their targets, however, were
not usually terrorists, but foreigners offering to spy for the United States or C.I.A. employees
suspected of misdeeds.

Agency officials, led by Mr.~Tenet, sought interrogation advice from other countries. And,
fatefully, they contacted the military unit that runs the SERE training program, the Joint Personnel
Recovery Agency, which gives American pilots, special operations troops and others a sample of the
brutal interrogation methods they might face as prisoners of war. Mr.~Tenet declined to be
interviewed.

By late 2001, the agency had contracted with James E.~Mitchell, a psychologist with the SERE program
who had monitored many mock interrogations but had never conducted any real ones, according to
colleagues. He was known for his belief that a psychological concept called ``learned helplessness''
was crucial to successful interrogation.

Martin Seligman, a prominent professor of psychology at the University of Pennsylvania who had
developed the concept, said in an interview that he was puzzled by Dr.~Mitchell's notion that
learned helplessness was relevant to interrogation.

``I think helplessness would make someone more dependent, less defiant and more compliant,''
Dr.~Seligman said, ``but I do not think it would lead reliably to more truth-telling.''

Still, forceful and brainy, Dr.~Mitchell, who declined to comment for this article, became a
persuasive player in high-level agency discussions about the best way to interrogate Qaeda
prisoners. Eventually, along with another former SERE psychologist, Bruce Jessen, Dr.~Mitchell
helped persuade C.I.A. officials that Qaeda members were fundamentally different from the myriad
personalities the agency routinely dealt with.

``Jim believed that people of this ilk would confess for only one reason: sheer terror,'' said one
C.I.A. official who had discussed the matter with Dr.~Mitchell.

Overwhelmed with reports of potential threats and anguished that the agency had failed to stop the
Sept.~11 attacks, Mr.~Tenet and his top aides did not probe deeply into the prescription
Dr.~Mitchell so confidently presented: using the SERE tactics on Qaeda prisoners.

A little research on the origin of those methods would have given reason for doubt. Government
studies in the 1950s found that Chinese Communist interrogators had produced false confessions from
captured American pilots not with some kind of sinister ``brainwashing'' but with crude tactics:
shackling the Americans to force them to stand for hours, keeping them in cold cells, disrupting
their sleep and limiting access to food and hygiene.

``The Communists do not look upon these assaults as 'torture,' '' one 1956 study concluded. ``But
all of them produce great discomfort, and lead to serious disturbances of many bodily processes;
there is no reason to differentiate them from any other form of torture.''

Worse, the study found that under such abusive treatment, a prisoner became ``malleable and
suggestible, and in some instances he may confabulate.''

In late 2001, about a half-dozen SERE trainers, according to a report released Tuesday night by the
Senate Armed Services Committee, began raising stark warning about plans by both the military and
the C.I.A. to use the SERE methods in interrogations.

In December 2001, Lt.~Col.~Daniel J.~Baumgartner of the Air Force, who oversaw SERE training,
cautioned in one memo that physical pressure was ``less reliable'' than other interrogation methods,
could backfire by increasing a prisoner's resistance and would have an ``intolerable public and
political backlash when discovered.'' But his memo went to the Defense Department, not the C.I.A.

One former senior intelligence official who played an important role in approving the interrogation
methods said he had no idea of the origins and history of the SERE program when the C.I.A. started
it in 2002.

``The agency was counting on the Justice Department to fully explore all the factors contributing to
a judgment about legality, including the surrounding history and context,'' the official said.

But it was the C.I.A. that was proposing the methods, and John Yoo, the Justice Department official
who was the principal author of a secret August 2002 memorandum that authorized the interrogation
program, was mostly interested in making a case that the president's wartime powers allowed for the
harsh tactics.

A Persuasive Case

After the March 28, 2002, capture in Pakistan of the Qaeda operative Abu Zubaydah -- the C.I.A.'s
first big catch after Sept.~11 -- Mr.~Tenet told Ms.~Rice, then the national security adviser, he
wanted to discuss interrogation, several former officials said. At a series of small-group and
individual briefings attended by Mr.~Bush, Mr.~Cheney, Ms.~Rice and Attorney General John Ashcroft,
Mr.~Tenet and his deputy, John McLaughlin, laid out their case.

They made a persuasive duo, former officials who heard their pitch recalled. Mr.~Tenet, an
extroverted former Congressional staff member, was given to forceful language about the threat from
Al Qaeda, which he said might well have had operations under way involving biological, radiological
or even nuclear weapons. Mr.~McLaughlin, a career intelligence analyst, was low-key and cerebral,
and some White House officials said they found his support for the methods reassuring.

In the briefings, Mr.~Tenet said that after extensive research, the agency believed that only the
methods he described -- which he said had been used on thousands of American trainees -- could
extract the details of plots from hardened Qaeda fanatics.

``It was described as a program that was safe and necessary, that would be closely monitored by
medical personnel,'' a former senior official recalled. ``And it was very much in the context of the
threat streams that were just eye-popping at the time.''

Mr.~Tenet's descriptions of each proposed interrogation method was so clinical and specific that at
one briefing Mr.~Ashcroft objected, saying that cabinet officials should approve broad outlines of
important policies, not the fine details, according to someone present. The attorney general later
complained that he thought Mr.~Tenet was looking for cover in case controversy erupted, the person
said.

Ms.~Rice insisted that Mr.~Ashcroft not just pass along the conclusions of his Office of Legal
Counsel, where Mr.~Yoo worked, but give his personal assurance that the methods were legal under
domestic and international law. He did.

The C.I.A. then gave individual briefings to the secretary of defense, Donald H.~Rumsfeld, and the
secretary of state, Colin L.~Powell. Neither objected, several former officials said.

Mr.~Cheney, whose top legal adviser, David S.~Addington, was closely consulting with Mr.~Yoo about
legal justification, strongly endorsed the program. Mr.~Bush also gave his approval, though what
details were shared with him is not known.

With that, the C.I.A. had the full support of the White House to begin its harshest interrogations.
Mr.~Bush and Mr.~Cheney have never publicly second-guessed their decision. Though some former
officials expressed regret that such a momentous decision was made so quickly without vital
information or robust debate, none were willing to be quoted by name.

There was one more check on intelligence programs, one designed in the 1970s to make sure
independent observers kept an eye on spy agencies: Congress. The Senate and House Intelligence
Committees had been created in the mid-1970s to prevent any repeat of the C.I.A. abuses unearthed by
the Senate's Church Committee.

As was common with the most secret programs, the C.I.A. chose not to brief the entire committees
about the interrogation methods but only the so-called Gang of Four -- the top Republican and
Democrat on the Senate and House committees. The rest of the committee members would be fully
briefed only in 2006.

The 2002 Gang of Four briefings left a hodgepodge of contradictory recollections that, to some
Congressional staff members, reveal a dysfunctional oversight system. Without full staff support,
few lawmakers are equipped to make difficult legal and policy judgments about secret programs,
critics say.

Representative Nancy Pelosi of California, who in 2002 was the ranking Democrat on the House
committee, has said in public statements that she recalls being briefed on the methods, including
waterboarding. She insists, however, that the lawmakers were told only that the C.I.A. believed the
methods were legal -- not that they were going to be used.

By contrast, the ranking Republican on the House committee at the time, Porter J.~Goss of Florida,
who later served as C.I.A. director, recalls a clear message that the methods would be used.

``We were briefed, and we certainly understood what C.I.A. was doing,'' Mr.~Goss said in an
interview. ``Not only was there no objection, there was actually concern about whether the agency
was doing enough.''

Senator Bob Graham, Democrat of Florida, who was committee chairman in 2002, said in an interview
that he did not recall ever being briefed on the methods, though government officials with access to
records say all four committee leaders received multiple briefings.

Senator Richard C.~Shelby of Alabama, the senior Republican on the committee, declined to discuss
the briefings.

Vicki Divoll, general counsel of the Senate Intelligence Committee in 2002 and a former C.I.A.
lawyer, would have been a logical choice to advise senators on the legal status of the interrogation
methods. But because of the restricted briefings, Ms.~Divoll learned about them only years later
from news media accounts.

Ms.~Divoll, who now teaches government at the United States Naval Academy, said the interrogation
issue revealed the perils of such restricted briefings.

``The very programs that are among the most risky and controversial, and that therefore should get
the greatest congressional oversight,'' she said, ``in fact get the least.''

\section{At Core of Detainee Fight: Did Methods Stop Attacks?}
WASHINGTON -- Even the most exacting truth commission may have a hard time determining for certain
whether brutal interrogations conducted by the Central Intelligence Agency helped keep the country
safe.

Last week's release of long-secret Justice Department interrogation memorandums has given rise to
starkly opposing narratives about what, if anything, was gained by the C.I.A.'s use of
waterboarding, wall-slamming and other physical pressure to shock and intimidate Qaeda operatives.

Senior Bush administration officials, led by Vice President Dick Cheney and cheered by many
Congressional Republicans, are fighting a rear-guard action in defense of their record. Only by
using the harshest methods, they insist, did the intelligence agency get the information it needed
to round up Qaeda killers and save thousands of American lives.

Even President Obama's new director of national intelligence, Dennis C.~Blair, wrote in a memorandum
to his staff last week that ``high value information came from interrogations in which these methods
were used,'' an assertion left out when the memorandum was edited for public release. By contrast,
Mr.~Obama and most of his top aides have argued that the use of those methods betrayed American
values -- and anyway, produced unreliable information. Those are a convenient pair of opinions, of
course: the moral balancing would be far trickier if the C.I.A. methods were demonstrated to have
been crucial in disrupting major plots.

For both sides, the political stakes are high, as proposals for a national commission to unravel the
interrogation story appear to be gaining momentum. Mr.~Obama and his allies need to discredit the
techniques he has banned. Otherwise, in the event of a future terrorist attack, critics may blame
his decision to rein in C.I.A. interrogators.

But if a strong case emerges that the Bush administration authorized torture and got nothing but
prisoners' desperate fabrications in return, that will tarnish what Mr.~Bush and Mr.~Cheney have
claimed as their greatest achievement: preventing new attacks after Sept.~11, 2001.

Within the agency, the necessity, effectiveness and legality of the interrogation methods have been
repeatedly subject to review. The agency's inspector general, John L.~Helgerson, studied the program
in 2004 and raised serious questions. According to former intelligence officials, that led to
separate reviews by an internal panel headed by Henry A.~Crumpton, a veteran counterterrorism
officer, and by two outsiders, Gardner Peckham, who had served as national security adviser to Newt
Gingrich, and John J.~Hamre, a former deputy defense secretary.

Their conclusions remain classified, but that could change now that the intelligence agency's
techniques have been made public. In a twist this week, Mr.~Cheney, a fierce defender of secrecy as
vice president, called for the release of more classified memorandums that he asserted prove the
effectiveness of the coercive techniques.

The second-guessing of the C.I.A.'s methods inside the government began long before Mr.~Obama's
election. The Federal Bureau of Investigation, the government agency with the greatest knowledge of
Al Qaeda in 2001, chose not to participate in the C.I.A. interrogation program after agents became
uneasy about the earliest use of harsh methods in 2002 on Abu Zubaydah, a long-sought terrorist
facilitator.

In an interview with Vanity Fair last year, the F.B.I. director since 2001, Robert S.~Mueller III,
was asked whether any attacks had been disrupted because of intelligence obtained through the
coercive methods. ``I don't believe that has been the case,'' Mr.~Mueller said. (A spokesman for
Mr.~Mueller, John Miller, said on Tuesday, ``The quote is accurate.'')

That assessment stands in sharp contrast to many assertions by Mr.~Bush and Mr.~Cheney, who on Fox
News on Sunday said of the methods: ``They did work. They kept us safe for seven years.''

Four successive C.I.A. directors have made similar claims, and the most recent, Michael V.~Hayden,
said in January that he believed the methods ``got the maximum amount of information'' from
prisoners, citing specifically Abu Zubaydah and Khalid Shaikh Mohammed, the chief 9/11 plotter.

Many intelligence officials, including some opposed to the brutal methods, confirm that the program
produced information of great value, including tips on early-stage schemes to attack tall buildings
on the West Coast and buildings in New York's financial district and Washington. Interrogation of
one Qaeda operative led to tips on finding others, until the leadership of the organization was
decimated. Removing from the scene such dedicated and skilled plotters as Mr.~Mohammed, or the
Indonesian terrorist known as Hambali, almost certainly prevented future attacks.

But which information came from which methods, and whether the same result might have been achieved
without the political, legal and moral cost of the torture controversy, is hotly disputed, even
inside the intelligence agency.

The Justice Department memorandums released last week illustrate how difficult it can be to assess
claims of effectiveness. One 2005 memorandum, for example, asserts that ``enhanced techniques'' used
on Abu Zubaydah and Mr.~Mohammed ``yielded critical information.''

But the memorandum then lists among Abu Zubaydah's revelations the identification of Mr.~Mohammed
and of an alleged radiological bomb plot by Jose Padilla, the American Qaeda associate. Both those
disclosures were made long before Abu Zubaydah was subjected to harsh treatment, according to
multiple accounts.

On Mr.~Mohammed, the record is murkier. The memorandum says that ``before the C.I.A. used enhanced
techniques,'' Mr.~Mohammed ``resisted giving any answers to questions about future attacks, 'Simply
noting, 'Soon, you will know.' ''

But the same memorandum reveals in a footnote that Mr.~Mohammed, captured on March 1, 2003, was
waterboarded 183 times that month. That striking number, which would average out to six
waterboardings a day, suggests that interrogators did not try a traditional, rapport-building
approach for long before escalating to their most extreme tool.

Mr.~Obama paid his first visit to the agency this week, and his reference to the interrogation issue
made for an awkward moment in which he sounded like a teacher gently correcting his pupils.

``Don't be discouraged that we have to acknowledge potentially we've made some mistakes,'' he said.
``That's how we learn.''

\section{U.S. Is Said to Push Chrysler to Prepare for Chapter 11}

DETROIT -- The Treasury Department is directing Chrysler to prepare a Chapter 11 bankruptcy filing
that could come as soon as next week, people with direct knowledge of the action said Thursday.

The Treasury has an agreement in principle with the United Automobile Workers union, whose members'
pensions and retiree health care benefits would be protected as a condition of the bankruptcy
filing, said these people, who asked for anonymity because they were not authorized to discuss the
case.

Moreover, Fiat of Italy would complete its alliance with Chrysler while the company is under
bankruptcy protection.

The only major question that remains unresolved is what happens to Chrysler's lenders, who hold
\$6.9 billion in company debt. The government's most recent offer, presented Wednesday, would give
the company's lenders about 22 cents on the dollar, or \$1.5 billion, and a 5 percent equity stake
in a reorganized Chrysler. Earlier this week, a steering committee of the lenders proposed that they
receive 65 cents on the dollar, or \$4.5 billion, and a 40 percent equity stake.

If no agreement is reached between the government and Chrysler's lenders, a nasty legal fight could
emerge in bankruptcy. The creditors' claims are backed by most of the company's collateral,
including plants, brands and equipment, and the senior lenders will argue that they have first claim
on those assets -- even over and above the government's debt.

``In a negotiation like this, everything is speculation until there's a deal,'' said an
administration official who did not want to be named because the talks are private. ``It should
surprise no one that the administration is planning on contingencies, but we remain focused on the
goal and engaged with all stakeholders to bring Chrysler and Fiat to a working partnership.''

A Treasury spokeswoman declined to comment.

A bankruptcy filing by Chrysler would be the first among Detroit's troubled automakers, who have
been mired in a devastating sales slump since last fall. Treasury is also working with General
Motors to prepare a possible bankruptcy case, and the terms of a Chrysler filing might offer a
glimpse into the shape of G.M.'s own filing.

Some analysts questioned whether the Treasury's steps to prepare a bankruptcy case were an effort to
put more pressure on lenders, with which it has exchanged proposals meant to reduce Chrysler's debt.
Chrysler faces an April 30 deadline from the Treasury, while G.M. faces a June 1 deadline in its own
efforts to draft a new restructuring plan.

Under the most likely assumptions, Treasury will provide the financing that Chrysler needs to
operate while under bankruptcy protection. The Canadian government is also expected to participate
in backing the company.

The Globe and Mail of Toronto reported the Canadian government's role on Thursday.

Last month, the Obama administration told Chrysler it would provide up to \$6 billion in financing
if Chrysler and Fiat could complete a deal by the end of this month. Fiat originally agreed to take
35 percent of Chrysler, but the stake was subsequently reduced to 20 percent. The administration
said it would provide up to \$6 billion in financing if the two companies agreed, on top of \$4
billion in federal assistance that Chrysler has already received.

Although the two companies have been holding discussions on an out-of-court agreement, a bankruptcy
case would allow Fiat to more easily select the assets of Chrysler that it wants to preserve, such
as dealerships, factories and the company's product development operations, these people said. The
approach, which relies upon Section 363 of the federal bankruptcy code, is somewhat similar to what
the government is planning in the case of G,M..

Then, Chrysler could sell or jettison any assets it does not want to keep, and cancel franchise
agreements with superfluous car dealers.

``Chrysler has consistently said that its viability will be enhanced through an alliance with Fiat,
as it represents a change in the company's business model that expands its global competitiveness,''
said Lori McTavish, Chrysler's vice president of communications. ``As we move forward in this
process, we believe it's important to keep all options open.''

Ms.~McTavish said Chrysler would continue to work through the end of the month, based on direction
given by the President's auto task force, ``to secure the support of the necessary stakeholders and
reach a successful conclusion that the administration and U.S.~Treasury deems appropriate.''

The U.A.W., Chrysler and Treasury have reached agreements in principle that would protect workers'
benefits, people with knowledge of the negotiations said, and a similar agreement is expected to be
reached as soon as this weekend with the Canadian Auto Workers union.

Once Chrysler emerges from bankruptcy protection, it would largely be owned by Fiat, the U.A.W., the
Treasury and its lenders, these people said. A bankruptcy filing would likely wipe out existing
equity stakeholders, notably Cerberus Capital Management, which took over the carmaker from Daimler
in 2007.

Ron Gettelfinger, the U.A.W.'s president, issued a statement on Wednesday saying that the union was
``continuing to work toward an agreement that will be in the best interest of Chrysler workers,
retirees and the communities where the company does business.''

People close to the talks said Wednesday that the U.A.W. had tentatively agreed to accept Chrysler
stock to finance half of the company's \$10.6 billion obligation to the health care trust. The
balance would be paid in cash over the next decade. That money presumably could come from either the
Treasury, or from Chrysler's profits, once it emerges from bankruptcy protection.

Chrysler has a \$9.3 billion pension shortfall, or 34 percent of its total liability, according to
the Pension Benefit Guaranty Corporation. The agency said earlier this month that it would assume
\$2 billion of the shortfall in the event Chrysler terminates its pension plans.

If that happened, retirees would receive sharply lower benefits than they normally would expect. But
Chrysler is not obligated to terminate its pension plans while in bankruptcy, particularly if it
received federal assistance to fund them.

It was not clear Thursday where Chrysler would file its bankruptcy case. On Wednesday, Mike Cox, the
attorney general of Michigan, urged General Motors and Chrysler to consider filing in the state,
rather than Delaware or New York. He said a locally administered case would be more convenient for
creditors in Michigan.

\section{Geithner, as Member and Overseer, Forged Ties to Finance Club}

Last June, with a financial hurricane gathering force, Treasury Secretary Henry M.~Paulson
Jr.~convened the nation's economic stewards for a brainstorming session. What emergency powers might
the government want at its disposal to confront the crisis? he asked.

Timothy F.~Geithner, who as president of the New York Federal Reserve Bank oversaw many of the
nation's most powerful financial institutions, stunned the group with the audacity of his answer. He
proposed asking Congress to give the president broad power to guarantee all the debt in the banking
system, according to two participants, including Michele Davis, then an assistant Treasury
secretary.

The proposal quickly died amid protests that it was politically untenable because it could put
taxpayers on the hook for trillions of dollars.

``People thought, 'Wow, that's kind of out there,' '' said John C.~Dugan, the comptroller of the
currency, who heard about the idea afterward. Mr.~Geithner says, ``I don't remember a serious
discussion on that proposal then.''

But in the 10 months since then, the government has in many ways embraced his blue-sky prescription.
Step by step, through an array of new programs, the Federal Reserve and Treasury have assumed an
unprecedented role in the banking system, using unprecedented amounts of taxpayer money, to try to
save the nation's financiers from their own mistakes.

And more often than not, Mr.~Geithner has been a leading architect of those bailouts, the activist
at the head of the pack. He was the federal regulator most willing to ``push the envelope,'' said
H.~Rodgin Cohen, a prominent Wall Street lawyer who spoke frequently with Mr.~Geithner.

Today, Mr.~Geithner is Treasury secretary, and as he seeks to rebuild the nation's fractured
financial system with more taxpayer assistance and a regulatory overhaul, he finds himself a locus
of discontent.

Even as banks complain that the government has attached too many intrusive strings to its financial
assistance, a range of critics -- lawmakers, economists and even former Federal Reserve colleagues
-- say that the bailout Mr.~Geithner has played such a central role in fashioning is overly generous
to the financial industry at taxpayer expense.

An examination of Mr.~Geithner's five years as president of the New York Fed, an era of unbridled
and ultimately disastrous risk-taking by the financial industry, shows that he forged unusually
close relationships with executives of Wall Street's giant financial institutions.

His actions, as a regulator and later a bailout king, often aligned with the industry's interests
and desires, according to interviews with financiers, regulators and analysts and a review of
Federal Reserve records.

In a pair of recent interviews and an exchange of e-mail messages, Mr.~Geithner defended his record,
saying that from very early on, he was ``a consistently dark voice about the potential risks ahead,
and a principal source of initiatives designed to make the system stronger'' before the markets
started to collapse.

Mr.~Geithner said his actions in the bailout were motivated solely by a desire to help businesses
and consumers. But in a financial crisis, he added, ``the government has to take risk, and we are
going to be doing things which ultimately -- in order to get the credit flowing again -- are going
to benefit the institutions that are at the core of the problem.''

The New York Fed is, by custom and design, clubby and opaque. It is charged with curbing banks'
risky impulses, yet its president is selected by and reports to a board dominated by the chief
executives of some of those same banks. Traditionally, the New York Fed president's
intelligence-gathering role has involved routine consultation with financiers, though Mr.~Geithner's
recent predecessors generally did not meet with them unless senior aides were also present,
according to the bank's former general counsel.

By those standards, Mr.~Geithner's reliance on bankers, hedge fund managers and others to assess the
market's health -- and provide guidance once it faltered -- stood out.

His calendars from 2007 and 2008 show that those interactions were a mix of the professional and the
private.

He ate lunch with senior executives from Citigroup, Goldman Sachs and Morgan Stanley at the Four
Seasons restaurant or in their corporate dining rooms. He attended casual dinners at the homes of
executives like Jamie Dimon, a member of the New York Fed board and the chief of JPMorgan Chase.

Mr.~Geithner was particularly close to executives of Citigroup, the largest bank under his
supervision. Robert E.~Rubin, a senior Citi executive and a former Treasury secretary, was
Mr.~Geithner's mentor from his years in the Clinton administration, and the two kept in close touch
in New York.

Mr.~Geithner met frequently with Sanford I.~Weill, one of Citi's largest individual shareholders and
its former chairman, serving on the board of a charity Mr.~Weill led. As the bank was entering a
financial tailspin, Mr.~Weill approached Mr.~Geithner about taking over as Citi's chief executive.

But for all his ties to Citi, Mr.~Geithner repeatedly missed or overlooked signs that the bank --
along with the rest of the financial system -- was falling apart. When he did spot trouble, analysts
say, his responses were too measured, or too late.

In 2005, for instance, Mr.~Geithner raised questions about how well Wall Street was tracking its
trading of complex financial products known as derivatives, yet he pressed reforms only at the
margins. Problems with the risky and opaque derivatives market later amplified the economic crisis.

As late as 2007, Mr.~Geithner advocated measures that government studies said would have allowed
banks to lower their reserves. When the crisis hit, banks were vulnerable because their financial
cushion was too thin to protect against large losses.

In fashioning the bailout, his drive to use taxpayer money to backstop faltering firms overrode
concerns that such a strategy would encourage more risk-taking in the future. In one bailout
instance, Mr.~Geithner fought a proposal to levy fees on banks that would help protect taxpayers
against losses.

The bailout has left the Fed holding a vast portfolio of troubled securities. To manage them,
Mr.~Geithner gave three no-bid contracts to BlackRock, an asset-management firm with deep ties to
the New York Fed.

To Joseph E.~Stiglitz, a Nobel-winning economist at Columbia and a critic of the bailout,
Mr.~Geithner's actions suggest that he came to share Wall Street's regulatory philosophy and world
view.

``I don't think that Tim Geithner was motivated by anything other than concern to get the financial
system working again,'' Mr.~Stiglitz said. ``But I think that mindsets can be shaped by people you
associate with, and you come to think that what's good for Wall Street is good for America.''

In this case, he added, that ``led to a bailout that was designed to try to get a lot of money to
Wall Street, to share the largesse with other market participants, but that had deeply obvious flaws
in that it put at risk the American taxpayer unnecessarily.''

But Ben S.~Bernanke, the chairman of the Federal Reserve, said in an interview that Mr.~Geithner's
Wall Street relationships made him ``invaluable'' as they worked together to steer the country
through crisis.

``He spoke frequently to many, many different players and kept his finger on the pulse of the
situation,'' Mr.~Bernanke said. ``He was the point person for me in many cases and with many
individual firms so that we were prepared for any kind of emergency.''

An Alternate Path

A revolving door has long connected Wall Street and the New York Fed. Mr.~Geithner's predecessors,
E.~Gerald Corrigan and William J.~McDonough, wound up as investment-bank executives. The current
president, William C.~Dudley, came from Goldman Sachs.

Mr.~Geithner followed a different route. An expert in international finance, he served under both
Clinton-era Treasury secretaries, Mr.~Rubin and Lawrence H.~Summers. He impressed them with his
handling of foreign financial crises in the late 1990s before landing a top job at the International
Monetary Fund.

When the New York Fed was looking for a new president, both former secretaries were advisers to the
bank's search committee and supported Mr.~Geithner's candidacy. Mr.~Rubin's seal of approval carried
particular weight because he was by then a senior official at Citigroup.

Mr.~Weill, Citigroup's architect, was a member of the New York Fed board when Mr.~Geithner arrived.
``He had a baby face,'' Mr.~Weill recalled. ``He didn't have a lot of experience in dealing with the
industry.''

But, he added, ``He quickly earned the respect of just about everyone I know. His knowledge, his
willingness to listen to people.''

At the age of 42, Mr.~Geithner took charge of a bank with enormous influence over the American
economy.

Sitting like a fortress in the heart of Manhattan's financial district, the New York Fed is, by dint
of the city's position as a world financial center, the most powerful of the 12 regional banks that
make up the Federal Reserve system.

The Federal Reserve was created after a banking crisis nearly a century ago to manage the money
supply through interest-rate policy, oversee the safety and soundness of the banking system and act
as lender of last resort in times of trouble. The Fed relies on its regional banks, like the New
York Fed, to carry out its policies and monitor certain banks in their areas.

The regional reserve banks are unusual entities. They are private and their shares are owned by
financial institutions the bank oversees. Their net income is paid to the Treasury.

At the New York Fed, top executives of global financial giants fill many seats on the board. In
recent years, board members have included the chief executives of Citigroup and JPMorgan Chase, as
well as top officials of Lehman Brothers and industrial companies like General Electric.

In theory, having financiers on the New York Fed's board should help the president be Washington's
eyes and ears on Wall Street. But critics, including some current and former Federal Reserve
officials, say the New York Fed is often more of a Wall Street mouthpiece than a cop.

Willem H.~Buiter, a professor at the London School of Economics and Political Science who caused a
stir at a Fed retreat last year with a paper concluding that the Federal Reserve had been co-opted
by the financial industry, said the structure ensured that ``Wall Street gets what it wants'' in its
New York president: ``A safe pair of hands, someone who is bright, intelligent, hard-working, but
not someone who intends to reform the system root and branch.''

Mr.~Geithner took office during one of the headiest bull markets ever. Yet his most important task,
he said in an interview, was to prepare banks for ``the storm that we thought was going to come.''

In his first speech as president in March 2004, he advised bankers to ``build a sufficient cushion
against adversity.'' Early on, he also spoke frequently about the risk posed by the explosion of
derivatives, unregulated insurancelike products that many companies use to hedge their bets.

But Mr.~Geithner acknowledges that ``even with all the things that we took the initiative to do, I
didn't think we achieved enough.''

Derivatives were not an altogether new issue for him, since the Clinton Treasury Department had
battled efforts to regulate the multitrillion-dollar market. As Mr.~Geithner shaped his own
approach, records and interviews show, he consulted veterans of that fight at Treasury, including
Lewis A.~Sachs, a close friend and tennis partner who managed a hedge fund.

Mr.~Geithner pushed the industry to keep better records of derivative deals, a measure that experts
credit with mitigating the chaos once firms began to topple. But he stopped short of pressing for
comprehensive regulation and disclosure of derivatives trading and even publicly endorsed their
potential to damp risk.

Nouriel Roubini, a professor of economics at the Stern School of Business at New York University,
who made early predictions of the crisis, said Mr.~Geithner deserved credit for trying, especially
given that the Fed chairman at the time, Alan Greenspan, was singing the praises of derivatives.

Even as Mr.~Geithner was counseling banks to take precautions against adversity, some economists
were arguing that easy credit was feeding a more obvious problem: a housing bubble.

Despite those warnings, a report released by the New York Fed in 2004 called predictions of gloom
``flawed'' and ``unpersuasive.'' And as lending standards evaporated and the housing boom reached
full throttle, banks plunged ever deeper into risky mortgage-backed securities and derivatives.

The nitty-gritty task of monitoring such risk-taking is done by 25 examiners at each large bank.
Mr.~Geithner reviewed his examiners' reports, but since they are not public, it is hard to fully
assess the New York Fed's actions during that period.

Mr.~Geithner said many of the New York Fed's supervisory actions could not be disclosed because of
confidentiality issues. As a result, he added, ``I realize I am vulnerable to a different narrative
in that context.''

The ultimate tool at Mr.~Geithner's disposal for reining in unsafe practices was to recommend that
the Board of Governors of the Fed publicly rebuke a bank with penalties or cease and desist orders.
Under his watch, only three such actions were taken against big domestic banks; none came after
2006, when banks' lending practices were at their worst.

The Citigroup Challenge

Perhaps the central regulatory challenge for Mr.~Geithner was Citigroup.

Cobbled together by Mr.~Weill through a series of pell-mell acquisitions into the world's largest
bank, Citigroup reached into every corner of the financial world: credit cards, auto loans, trading,
investment banking, as well as mortgage securities and derivatives. But it was plagued by
mismanagement and wayward banking practices.

In 2004, the New York Fed levied a \$70 million penalty against Citigroup over the bank's lending
practices. The next year, the New York Fed barred Citigroup from further acquisitions after the bank
was involved in trading irregularities and questions about its operations. The New York Fed lifted
that restriction in 2006, citing the company's ``significant progress'' in carrying out risk-control
measures.

In fact, risk was rising to dangerous levels at Citigroup as the bank dove deeper into
mortgage-backed securities.

Throughout the spring and summer of 2007, as subprime lenders began to fail and government officials
reassured the public that the problems were contained, Mr.~Geithner met repeatedly with members of
Citigroup's management, records show.

From mid-May to mid-June alone, he met over breakfast with Charles O.~Prince, the company's chief
executive at the time, traveled to Citigroup headquarters in Midtown Manhattan to meet with Lewis
B.~Kaden, the company's vice chairman, and had coffee with Thomas G.~Maheras, who ran some of the
bank's biggest trading operations.

(Mr.~Maheras's unit would later be roundly criticized for taking many of the risks that led
Citigroup aground.)

His calendar shows that during that period he also had breakfast with Mr.~Rubin. But in his
conversations with Mr.~Rubin, Mr.~Geithner said, he did not discuss bank matters. ``I did not do
supervision with Bob Rubin,'' he said.

Any intelligence Mr.~Geithner gathered in his meetings does not appear to have prepared him for the
severity of the problems at Citigroup and beyond.

In a May 15, 2007, speech to the Federal Reserve Bank of Atlanta, Mr.~Geithner praised the strength
of the nation's top financial institutions, saying that innovations like derivatives had ``improved
the capacity to measure and manage risk'' and declaring that ``the larger global financial
institutions are generally stronger in terms of capital relative to risk.''

Two days later, interviews and records show, he lobbied behind the scenes for a plan that a
government study said could lead banks to reduce the amount of capital they kept on hand.

While waiting for a breakfast meeting with Mr.~Weill at the Four Seasons Hotel in Manhattan,
Mr.~Geithner phoned Mr.~Dugan, the comptroller of the currency, according to both men's calendars.
Both Citigroup and JPMorgan Chase were pushing for the new standards, which they said would make
them more competitive. Records show that earlier that week, Mr.~Geithner had discussed the issue
with JPMorgan's chief, Mr.~Dimon.

At the Federal Deposit Insurance Corporation, which insures bank deposits, the chairwoman, Sheila
C.~Bair, argued that the new standards were tantamount to letting the banks set their own capital
levels. Taxpayers, she warned, could be left ``holding the bag'' in a downturn. But Mr.~Geithner
believed that the standards would make the banks more sensitive to risk, Mr.~Dugan recalled. The
standards were adopted but have yet to go into effect.

Callum McCarthy, a former top British financial regulator, said regulators worldwide should have
focused instead on how undercapitalized banks already were. ``The problem is that people in banks
overestimated their ability to manage risk, and we believed them.''

By the fall of 2007, that was becoming clear. Citigroup alone would eventually require \$45 billion
in direct taxpayer assistance to stay afloat.

On Nov. 5, 2007, Mr.~Prince stepped down as Citigroup's chief in the wake of multibillion-dollar
mortgage write-downs. Mr.~Rubin was named chairman, and the search for a new chief executive began.
Mr.~Weill had a perfect candidate: Mr.~Geithner.

The two men had remained close. That past January, Mr.~Geithner had joined the board of the National
Academy Foundation, a nonprofit organization founded by Mr.~Weill to help inner-city high school
students prepare for the work force.

``I was a little worried about the implications,'' Mr.~Geithner said, but added that he had accepted
the unpaid post only after Mr.~Weill had stepped down as Citigroup's chairman, and because it was a
good cause that the Fed already supported.

Although Mr.~Geithner was a headliner with Mr.~Prince at a 2004 fundraiser that generated \$1.1
million for the foundation, he said he did not raise money for the group once on the board. He
attended regular foundation meetings at Mr.~Weill's Midtown Manhattan office.

In addition to charity business, Mr.~Weill said, the two men often spoke about what was happening at
Citigroup. ``It would be logical,'' he said.

On Nov. 6 and 7, 2007, as Mr.~Geithner's bank examiners scrambled to assess Citigroup's problems,
the two men spoke twice, records show, once for a half-hour on the phone and once for an hourlong
meeting in Mr.~Weill's office, followed by a National Academy Foundation cocktail reception.

Mr.~Geithner also went to Citigroup headquarters for a lunch with Mr.~Rubin on Nov. 16 and met with
Mr.~Prince on Dec.~4, records show.

Mr.~Geithner acknowledged in an interview that Mr.~Weill had spoken with him about the Citigroup
job. But he immediately rejected the idea, he said, because he did not think he was right for the
job.

``I told him I was not the right choice,'' Mr.~Geithner said, adding that he then spoke to ``one
other board member to confirm after the fact that it did not make sense.''

According to New York Fed officials, Mr.~Geithner informed the reserve bank's lawyers about the
exchange with Mr.~Weill, and they told him to recuse himself from Citigroup business until the
matter was resolved.

Mr.~Geithner said he ``would never put myself in a position where my actions were influenced by a
personal relationship.''

Other chief financial regulators at the Federal Deposit Insurance Company and the Securities and
Exchange Commission say they keep officials from institutions they supervise at arm's length, to
avoid even the appearance of a conflict. While the New York Fed's rules do not prevent its president
from holding such one-on-one meetings, that was not the general practice of Mr.~Geithner's recent
predecessors, said Ernest T.~Patrikis, a former general counsel and chief operating officer at the
New York Fed.

``Typically, there would be senior staff there to protect against disputes in the future as to the
nature of the conversations,'' he said.

Coping With Crisis

As Mr.~Geithner sees it, most of the institutions hit hardest by the crisis were not under his
jurisdiction -- some foreign banks, mortgage companies and brokerage firms. But he acknowledges that
``the thing I feel somewhat burdened by is that I didn't attempt to try to change the rules of the
game on capital requirements early on,'' which could have left banks in better shape to weather the
storm.

By last fall, it was too late. The government, with Mr.~Geithner playing a lead role alongside
Mr.~Bernanke and Mr.~Paulson, scurried to rescue the financial system from collapse. As the Fed
became the biggest vehicle for the bailout, its balance sheet more than doubled, from \$900 billion
in October 2007 to more than \$2 trillion today.

``I couldn't have cared less about Wall Street, but we faced a crisis that was going to cause
enormous damage to the economy,'' Mr.~Geithner said.

The first to fall was Bear Stearns, which had bet heavily on mortgages and by mid-March was
tottering. Mr.~Geithner and Mr.~Paulson persuaded JPMorgan Chase to take over Bear. But to complete
the deal, JPMorgan insisted that the government buy \$29 billion in risky securities owned by Bear.

Some officials at the Federal Reserve feared encouraging risky behavior by bailing out an investment
house that did not even fall under its umbrella. To Mr.~Geithner's supporters, that he prevailed in
the case of Bear and other bailout decisions is testament to his leadership.

``He was a leader in trying to come up with an aggressive set of policies so that it wouldn't get
completely out of control,'' said Philipp Hildebrand, a top official at the Swiss National Bank who
has worked with Mr.~Geithner to coordinate an international response to the worldwide financial
crisis.

But others are less enthusiastic. William Poole, president of the Federal Reserve Bank of St. Louis
until March 2008, said that the Fed, by effectively creating money out of thin air, not only runs
the risk of ``massive inflation'' but has also done an end-run around Congressional power to control
spending.

Many of the programs ``ought to be legislated and shouldn't be in the Federal Reserve at all,'' he
contended.

In making the Bear deal, the New York Fed agreed to accept Bear's own calculation of the value of
assets acquired with taxpayer money, even though those values were almost certain to decline as the
economy deteriorated. Although Fed officials argue that they can hold onto those assets until they
increase in value, to date taxpayers have lost \$3.4 billion. Even these losses are probably
understated, given how the Federal Reserve priced the holdings, said Janet Tavakoli, president of
Tavakoli Structured Finance, a consulting firm in Chicago. ``You can assume that it has used magical
thinking in valuing these assets,'' she said.

Mr.~Geithner played a pivotal role in the next bailout, which was even bigger -- that of the
American International Group, the insurance giant whose derivatives business had brought it to the
brink of collapse in September. He also went to bat for Goldman Sachs, one of the insurer's biggest
trading partners.

As A.I.G. bordered on bankruptcy, Mr.~Geithner pressed first for a private sector solution. A.I.G.
needed \$60 billion to meet payments on insurance contracts it had written to protect customers
against debt defaults.

A.I.G.'s chief executive at the time, Robert B.~Willumstad, said he had hired bankers at JPMorgan to
help it raise capital. Goldman Sachs had jockeyed for the job as well, but because the investment
bank was one of A.I.G.'s biggest trading partners, Mr.~Willumstad rejected the idea. The potential
conflicts of interest, he believed, were too great.

Nevertheless, on Monday, Sept.~15, Mr.~Geithner pushed A.I.G. to bring Goldman onto its team to
raise capital, Mr.~Willumstad said.

Mr.~Geithner and Mr.~Corrigan, a Goldman managing director, were close, speaking frequently and
sometimes lunching together at Goldman headquarters. On that day, the company's chief executive,
Lloyd C.~Blankfein, was at the New York Fed.

A Goldman spokesman said, ``We don't believe anyone at Goldman Sachs asked Mr.~Geithner to include
the firm in the assignment.'' Mr.~Geithner said he had suggested Goldman get involved because the
situation was chaotic and ``time was running out.''

But A.I.G.'s search for capital was fruitless. By late Tuesday afternoon, the government would step
in with an \$85 billion loan, the first installment of a bailout that now stands at \$182 billion.
As part of the bailout, A.I.G.'s trading partners, including Goldman, were compensated fully for
money owed to them by A.I.G.

Analysts say the New York Fed should have pressed A.I.G.'s trading partners to take a deep discount
on what they were owed. But Mr.~Geithner said he had no bargaining power because he was unwilling to
threaten A.I.G.'s trading partners with a bankruptcy by the insurer for fear of further
destabilizing the system.

A recent report on the A.I.G. bailout by the Government Accountability Office found that taxpayers
may never get their money back.

The Debt Guarantee

Over Columbus Day weekend last fall, with the market gripped by fear and banks refusing to lend to
one another, a somber group gathered in an ornate conference room across from Mr.~Paulson's office
at the Treasury.

Mr.~Paulson, Mr.~Bernanke, Ms.~Bair and others listened as Mr.~Geithner made his pitch, according to
four participants. Mr.~Geithner, in the words of one participant, was ``hell bent'' on a plan to use
the Federal Deposit Insurance Corporation to guarantee debt issued by bank holding companies.

It was a variation on Mr.~Geithner's once-unthinkable plan to have the government guarantee all bank
debt.

The idea of putting the government behind debt issued by banking and investment companies was a
momentous shift, an assistant Treasury secretary, David G.~Nason, argued. Mr.~Geithner wanted to
give the banks the guarantee free, saying in a recent interview that he felt that charging them
would be ``counterproductive.'' But Ms.~Bair worried that her agency -- and ultimately taxpayers --
would be left vulnerable in the event of a default.

Mr.~Geithner's program was enacted and to date has guaranteed \$340 billion in loans to banks. But
Ms.~Bair prevailed on taking fees for the guarantees, and the government so far has collected \$7
billion.

Mr.~Geithner has also faced scrutiny over how well taxpayers were served by his handling of another
aspect of the bailout: three no-bid contracts the New York Fed awarded to BlackRock, a money
management firm, to oversee troubled assets acquired by the bank.

BlackRock was well known to the Fed. Mr.~Geithner socialized with Ralph L.~Schlosstein, who founded
the company and remains a large shareholder, and has dined at his Manhattan home. Peter R.~Fisher,
who was a senior official at the New York Fed until 2001, is a managing director at BlackRock.

Mr.~Schlosstein said that while he and Mr.~Geithner spoke frequently, BlackRock's work for the Fed
never came up.

``Conversations with Tim were appropriately a one-way street. He'd call you and pepper you with a
bunch of questions and say thank you very much and hang up,'' he said. ``My experience with Tim is
that he makes those kinds of decisions 100 percent based on capability and zero about
relationships.''

For months, New York Fed officials declined to make public details of the contract, which has become
a flash point with some lawmakers who say the Fed's handling of the bailout is too secretive. New
York Fed officials initially said in interviews that they could not disclose the fees because they
had agreed with BlackRock to keep them confidential in exchange for a discount.

The contract terms they subsequently disclosed to The New York Times show that the contract is worth
at least \$71.3 million over three years. While that rate is largely in keeping with comparable fees
for such services, analysts say it is hardly discounted.

Mr.~Geithner said he hired BlackRock because he needed its expertise during the Bear
Stearns-JPMorgan negotiations. He said most of the other likely candidates had conflicts, and he had
little time to shop around. Indeed, the deal was cut so quickly that they worked out the fees only
after the firm was hired.

But since then, the New York Fed has given two more no-bid contracts to BlackRock related to the
A.I.G. bailout, angering a number of BlackRock's competitors. The fees on those contracts remain
confidential.

Rescues Revisited

As Mr.~Geithner runs the Treasury and administration officials signal more bailout money may be
needed, the specter of bailouts past haunts his efforts.

He recently weathered a firestorm over retention payments to A.I.G. executives made possible in part
by language inserted in the administration's stimulus package at the Treasury Department's
insistence. And his new efforts to restart the financial industry suggest the same philosophy that
guided Mr.~Geithner's Fed years.

According to a recent report by the inspector general monitoring the bailout, Neil M.~Barofsky,
Mr.~Geithner's plan to underwrite investors willing to buy the risky mortgage-backed securities
still weighing down banks' books is a boon for private equity and hedge funds but exposes taxpayers
to ``potential unfairness'' by shifting the burden to them.

The top echelon of the Treasury Department is a common destination for financiers, and Mr.~Geithner
has also recruited aides from Wall Street, some from firms that were at the heart of the crisis. For
instance, his chief of staff, Mark A.~Patterson, is a former lobbyist for Goldman Sachs, and one of
his top counselors is Lewis S.~Alexander, a former chief economist at Citigroup.

A bill sent recently by the Treasury to Capitol Hill would give the Obama administration extensive
new powers to inject money into or seize systemically important firms in danger of failure. It was
drafted in large measure by Davis Polk \& Wardwell, a law firm that represents many banks and the
financial industry's lobbying group. Mr.~Geithner also hired Davis Polk to represent the New York
Fed during the A.I.G. bailout.

Treasury officials say they inadvertently used a copy of Davis Polk's draft sent to them by the
Federal Reserve as a template for their own bill, with the result that the proposed legislation
Treasury sent to Capitol Hill bore the law firm's computer footprints. And they point to several
significant changes to that draft that ``better protect the taxpayer,'' in the words of Andrew
Williams, a Treasury spokesman.

But others say important provisions in the original industry bill remain. Most significant, the bill
does not require that any government rescue of a troubled firm be done at the lowest possible cost,
as is required by the F.D.I.C. when it takes over a failed bank. Treasury officials said that is
because they would use the rescue powers only in rare and extreme cases that might require
flexibility. Karen Shaw Petrou, managing director of the Washington research firm Federal Financial
Analytics, said it essentially gives Treasury ``a blank check.''

One year and two administrations into the bailout, Mr.~Geithner is perhaps the single person most
identified with the enormous checks the government has written. At every turn, he is being
second-guessed about the rescues' costs and results. But he remains firm in his belief that failure
to act would have been much more costly.

``All financial crises are a fight over how much losses the government ultimately takes on,'' he
said. And every decision ``requires we balance how to achieve the most benefits in terms of
improving confidence and the flow of credit at the least risk to taxpayers.''

\section{China Still Presses Crusade Against Falun Gong}

BEIJING -- In the decade since the Chinese government began repressing Falun Gong, a crusade that
human rights groups say has led to the imprisonment of tens of thousands of practitioners and
claimed at least 2,000 lives, the world's attention has shifted elsewhere.

The drive against the spiritual group has eliminated its leadership, decimated the ranks of faithful
and convinced many Chinese that the group is an ``evil cult,'' as the government contends. But 10
years on, the war on Falun Gong remains unfinished.

In the past year, as many as 8,000 practitioners have been detained, according to experts on human
rights, and at least 100 have died in custody. Among them were Yu Zhou, 42, a popular Beijing
musician, and Cao Changling, the 77-year-old vice director of a paper plant in Wuhan, whose bruised
body was returned to his family by the police last summer just as China was reveling in the glory of
the Olympic Games.

In recent months, scores of practitioners have been given long prison terms, including Zhang Xingwu,
a retired physics professor from Shandong Province who last week was sentenced to seven years after
the police found Falun Gong literature in his apartment, according to family members.

The continued crackdown highlights the difficulty of eradicating a movement whose adherents
stubbornly cling to their beliefs, but it also provides a window into the psyche of an authoritarian
government that, despite its far-reaching power, remains deeply insecure.

From the outset, the group, which at its peak claimed to have millions of followers around China,
insisted that it wanted only legal recognition, not political power. But the country's top leaders
were alarmed by the group's ability to attract a devoted following from so many citizens -- from
retired functionaries to pimple-faced college students.

The decision to ban the group entirely was made after 10,000 Falun Gong adherents staged a silent
protest outside the gates of Zhongnanhai, the Communist Party's leadership compound in Beijing, to
complain about reports in the state-run media that the group said were defamatory. Security forces
apparently had no advance knowledge of the demonstration, which took place on April 25, 1999, and
they began treating the group as a threat to national security.

``Even a soccer team with an organization like Falun Gong might have produced the same reaction,''
said T.~Kumar, the Asia advocacy director for Amnesty International.

Although the propaganda juggernaut has eased in recent years, Falun Gong remains a toxic subject in
China. Few academics will speak about it on the record, and the Internet is scoured clean of
information that might be construed as sympathetic to Falun Gong, an amalgam of Buddhism, mysticism
and qigong, the traditional exercise regimen that remains broadly popular here.

For the Falun Gong devotees who practice in secret, the only glimmer of hope has come from a small
but growing number of lawyers who have dared to take on their cases. Even if the legal efforts have
mostly come to naught, until recently Falun Gong detainees were denied even the right to a lawyer.

Last week, Jiang Yu, a spokesman for the Foreign Ministry, reiterated the government's long-held
stance that Falun Gong warrants suppression because it emphasizes meditation and the paranormal over
modern medicine. ``The Falun Gong cult violates human rights by controlling people's minds,'' he
said in response to a reporter's query.

Among experts based outside the country, there is broad consensus that the government's efforts have
not done much to advance its own interests, at least internationally, where it has been dogged by
allegations that it uses torture to crush believers into submission.

'The excesses and the savagery have really lowered the quality of the government and harmed its
reputation abroad,'' said Jerome Cohen, a member of the Council on Foreign Relations and an expert
on Chinese law. ``They're paying a high price for the cruelty to these people.''

According to Falun Gong followers and Chinese lawyers who take on their cases, that cruelty
continues unabated.

Among those swept up in the purge were Yu Zhou, the musician, and his wife, Xu Na. They were stopped
for speeding in January 2008, according to their lawyer. After the police found Falun Gong materials
in their car, both were detained. Ten days later, Mr.~Yu's sister was told that he was gravely ill,
the result, she was told, of a hunger strike complicated by diabetes. His sister, Yu Qun, says her
brother did not have diabetes. She contends that he died at the hands of his captors.

The family's efforts to investigate Mr.~Yu's death have been thwarted by the police and prosecutors,
who refuse to allow an autopsy or even issue a death certificate.

Ms.~Xu, who is a well-known poet and painter, was given a three-year term.

``I don't understand why this happened to them because they didn't do anything to break the law and
they weren't promoting the group,'' Ms.~Yu said.

According to former detainees and human rights organizations, Falun Gong detainees are frequently
subjected to harrowing abuse, particularly those who refuse to swear off their faith. Bu Dongwei,
41, a longtime adherent who spent three years in a labor camp, said he was forced to share a room
with about 30 people, most of them petty thieves and drug addicts who were encouraged to abuse the
Falun Gong detainees.

Mr.~Bu, a trained geneticist, left China in December and now lives in Los Angeles.

While the group's initial goals were official legitimacy and an end to persecution, the ceaseless
campaign against them has radicalized many adherents, especially those living outside China. In
cities around the world, Falun Gong devotees -- and their offbeat re-enactments of torture and gory
visual aids -- have become a common sight. The group has dedicated itself to the demise of the
Communist Party, which has complicated the lives of adherents inside China.

Falun Dafa, the organization that oversees the movement from its headquarters in New York, is led by
Li Hongzhi, a former grain clerk who began spreading his mystical brand of qigong in 1992 but fled
China before the crackdown began. Once known for charismatic preaching, he has spent much of the
past decade living a reclusive life in Queens.

David Ownby, the author of ``Falun Gong and the Future of China,'' said that Mr.~Li and his
followers may have made a tactical mistake by massing in Beijing, but that the Communist Party erred
by interpreting their actions as a threat to its rule.

``If either side had played their cards more intelligently, Falun Gong could have been co-opted by
the government,'' said Mr.~Ownby, who is a professor of East Asian studies at the University of
Montreal. He added, ``This horrific loss of life could have been avoided.''

\section{In China, Knockoff Cellphones Are a Hit}

SHENZHEN, China -- The phone's sleek lines and touch-screen keyboard are unmistakably familiar. So
is the logo on the back. But a sales clerk at a sprawling electronic goods market in this Chinese
coastal city admits what is clear upon closer inspection: this is not the Apple iPhone; this is the
Hi-Phone.

``But it's just as good,'' the clerk says.

Nearby, dozens of other vendors are selling counterfeit Nokia, Motorola and Samsung phones -- as
well as cheap look-alikes that make no bones about being knockoffs.

``Five years ago, there were no counterfeit phones,'' says Xiong Ting, a sales manager at Triquint
Semiconductor, a maker of mobile phone parts, while visiting Shenzhen. ``You needed a design house.
You needed software guys. You needed hardware design. But now, a company with five guys can do it.
Within 100 miles of here, you can find all your suppliers.''

Technological advances have allowed hundreds of small Chinese companies, some with as few as 10
employees, to churn out what are known here as shanzhai, or black market, cellphones, often for as
little as \$20 apiece.

And just as Chinese companies are trying to move up the value chain of manufacturing, from producing
toys and garments to making computers and electric cars, so too are counterfeiters. After years of
making fake luxury bags and cheap DVDs, they are capturing market share from the world's biggest
mobile phone makers.

Although shanzhai phones have only been around a few years, they already account for more than 20
percent of sales in China, which is the world's biggest mobile phone market, according to the
research firm Gartner.

They are also being illegally exported to Russia, India, the Middle East, Europe, even the United
States. ``The shanzhai phone market is expanding crazily,'' says Wang Jiping, a senior analyst at
IDC, which tracks technology trends. ``They copy Apple, Nokia, whatever they like, and they respond
to the market swiftly.''

Alarmed by the rapid growth of counterfeits and no-name knockoffs, global brands are pressing the
Chinese government to crack down on their proliferation, and are warning consumers about potential
health hazards, like cheap batteries that can explode.

Nokia, the world's biggest cellphone maker, says it is working with Beijing to fight counterfeiting.
Motorola says much the same. Apple Inc.~declined to comment.

Even Chinese mobile phone producers are losing market share to underground companies, which have a
built-in cost advantage because they evade taxes, regulatory fees and safety checks.

``We're being severely hurt by shanzhai phones,'' says Chen Zhao, a sales director at Konka, a
Chinese cellphone maker. ``Legal cellphone makers should pay 17 percent of their revenue as
value-added tax, but shanzhai makers, of course, won't pay it.''

So far, however, China has done little to stop the proliferation of fake mobile phones, which are
even advertised on late-night television infomercials with pitches like ``one-fifth the price, but
the same function and look,'' or patriotic appeals like ``Buy shanzhai to show your love of our
country.''

Last month, the Ministry of Industry and Information Technology did warn consumers about the hazards
of shanzhai phones, saying ``their radiation usually exceeds the limit.'' China's consumer
protection agency says faulty mobile phones were the No.~1 consumer complaint last year.

A few weeks ago, a 45-year-old man in south China was severely burned after his cellphone exploded
in his shirt pocket, according to state-run news media.

But that hasn't seemed to affect sales of black market phones, which typically sell at retail for
\$100 to \$150. In the spirit of what is called ``shanzhai'' -- which suggests rebels or bandits and
which applies to counterfeit products of all kinds -- many consumers are willing to take a risk on a
cheap item that looks stylish.

``I saw iPhone pictures on the Web; it's so cool. But it costs over \$500 -- too expensive,'' says
Yang Guibin, 30, an office worker from Chongqing. ``So I decided to buy a shanzhai iPhone. I bought
it in a digital market here; it looked exactly like the iPhone.''

Some experts say they believe the shanzhai phenomena is about being creative, Chinese style.

``Chinese grass-roots companies are actually very innovative,'' says Yu Zhou, a professor at Vassar
College. ``It's not so much technology as how they form supply chains and how rapidly they react to
new trends.''

While the phones may look like famous brands, companies actually add special features like bigger
screens, dual-mode SIM cards (which allow two phone numbers) and even a telescopic lens attachment
for the phone's camera.

Since it is the SIM card that makes a phone run in China, as in most places other than the United
States, all you have to do is insert a valid SIM card into a shanzhai phone and it works.

All this innovation comes from an industry that only took off in 2005, after Mediatek, a
semiconductor design company from Taiwan, helped significantly reduce the cost and complexity of
producing a mobile phone.

Using what experts call a turnkey solution, Mediatek developed a circuit board that could
inexpensively integrate the functions of multiple chips, offering start-ups a platform to produce a
low-cost mobile phone.

The industry got another boost in 2007, when regulators said companies no longer needed a license to
manufacture a cellphone.

That set off a scramble by entrepreneurs in this electronics manufacturing center. Counterfeiting
and off-brand knockoffs flourished. Tiny companies would buy a Mediatek chip loaded with software,
source other components and ask a factory to assemble them.

Marketing strategies were simple: steal. Designs and brand names were copied identically or simply
mimicked. (Sumsung for Samsung or Nckia for Nokia.)

Tapping into the supply chains of big brands is easy, producers say. ``It's really common for
factories to do a night shift for other companies,'' says Zhang Haizhen, who recently ran a shanzhai
company here. ``No one will refuse an order if it is over 5,000 mobile phones.''

The people who make fake iPhones admit it's a shady business.

``We are a kind of illegal producer,'' says Zhang Feiyang, whose company, Yuanyang, makes an iPhone
clone. ``In Shenzhen there are many small mills, hidden. Basically, we can make any type of
cellphone.''

The competition is already forcing global brands to lower prices, analysts say. And new Chinese
brands are emerging, like Meizu, a would-be Apple that has opened stylish stores here.

``Our phone is even better than the iPhone,'' says Liu Zeyu, a Meizu salesman in Shenzhen. ``Our
goal is to create a phone that makes Chinese proud.''

\section{In Developing Countries, Web Grows Without Profit}

Facebook is booming in Turkey and Indonesia. YouTube's audience has nearly doubled in India and
Brazil.

That may seem like good news. But it is also a major reason these and other Web companies with big
global audiences and renowned brands struggle to turn even a tiny profit.

Call it the International Paradox.

Web companies that rely on advertising are enjoying some of their most vibrant growth in developing
countries. But those are also the same places where it can be the most expensive to operate, since
Web companies often need more servers to make content available to parts of the world with limited
bandwidth. And in those countries, online display advertising is least likely to translate into
results.

This intractable contradiction has become a serious drag on the bottom lines of photo-sharing sites,
social networks and video distributors like YouTube. It is also threatening the fervent idealism of
Internet entrepreneurs, who hoped to unite the world in a single online village but are increasingly
finding that the economics of that vision just do not work.

Last year, Veoh, a video-sharing site operated from San Diego, decided to block its service from
users in Africa, Asia, Latin America and Eastern Europe, citing the dim prospects of making money
and the high cost of delivering video there.

``I believe in free, open communications,'' Dmitry Shapiro, the company's chief executive, said.
``But these people are so hungry for this content. They sit and they watch and watch and watch. The
problem is they are eating up bandwidth, and it's very difficult to derive revenue from it.''

Internet start-ups that came of age during the Web 2.0 era, roughly from 2004 to the beginning of
the recession at the end of 2007, generally subscribed to a widely accepted blueprint: build huge
global audiences with a free service, and let advertising pay the bills.

But many of them ran smack into global economic reality. There may be 1.6 billion people in the
world with Internet access, but fewer than half of them have incomes high enough to interest major
advertisers.

``It's a problem every Internet company has,'' said Michelangelo Volpi, chief executive of Joost, a
video site with half its audience outside the United States.

``Whenever you have a lot of user-generated material, your bandwidth gets utilized in Asia, the
Middle East, Latin America, where bandwidth is expensive and ad rates are ridiculously low,'' Mr.
Volpi said. If Web companies ``really want to make money, they would shut off all those countries.''

Few Internet companies have taken that drastic step, but many are exploring other ways to increase
revenue or cut costs in developing countries.

MySpace -- the News Corporation's social network with 130 million members, about 45 percent of them
overseas -- is testing a feature for countries with slower Internet connections called Profile Lite.
It is a stripped-down version of the site that is less expensive to display because it requires less
bandwidth.

MySpace says it may make Profile Lite the primary version for its members in India, where it has
760,000 users, although people there could click on a link to switch to the richer version of the
site.

Perhaps no company is more in the grip of the international paradox than YouTube, which a Credit
Suisse analyst, Spencer Wang, recently estimated could lose \$470 million in 2009, in part because
of the high cost of delivering billions of videos each month. Google, which owns YouTube, disputed
the analysis but offered no details on the site's financial situation.

Tom Pickett, director of online sales and operations at YouTube, says the company still hews to its
vision of bringing online video to the entire globe. In the last two years, it has pushed to create
local versions of its site in countries like India, Brazil and Poland.

But Mr.~Pickett also says that YouTube has slowed the creation of new international hubs and shifted
its focus to making money. He says that does not rule out restricting bandwidth in certain countries
as a way to control costs -- essentially making YouTube a slower, lower-quality viewing experience
in the developing world.

``We may choose to set a limit to how much we are willing to pay in bandwidth cost,'' Mr.~Pickett
said. In some countries, he said, ``there may be particular peak times where instead of high
definition, we might decrease the resolution.''

The Facebook social network is also considering lowering the quality of videos and photographs
delivered to some regions in an effort to reduce expenses.

``We can decide, either on a country by country or user by user basis, to engineer the quality of
the service for that cohort of users,'' said Jonathan Heiliger, the executive who oversees
Facebook's computing infrastructure.

Facebook is in a particularly difficult predicament. Seventy percent of its 200 million members live
outside the United States, many in regions that do not contribute much to Facebook's bottom line. At
the same time, the company faces the expensive prospect of storing 850 million photos and eight
million videos uploaded to the site each month.

Facebook, which says it favors membership growth over profitability for now, is trying to increase
revenue overseas by hiring advertising sales staff in countries like Britain, Australia and France.

In other parts of the world, Microsoft serves ads on the site and Facebook offers self-service tools
to advertisers. But those ads are far less lucrative than the ones Facebook itself sells in the
United States and Western Europe.

As a result, speculation has swirled about Facebook's finances. Industry analysts wonder aloud how
fast the company is losing money and whether it needs to solicit another round of investment.

Facebook said last month that it was on track to become profitable next year. But as it did, Gideon
Yu, Facebook's experienced chief financial officer, left the company. Three people familiar with the
internal maneuverings at Facebook said Mr.~Yu objected to such a rosy projection as the company was
struggling to finance its expensive global growth.

Web entrepreneurs like Mr.~Shapiro of Veoh, still struggling with his decision to restrict his site
from much of the world, might have to find a way to soothe their battered consciences.

``The part of me that wants to change the world says, 'This is unfair, it shouldn't be like this,'
'' Mr.~Shapiro said. ``On the other hand, from the business side of things, serving videos to the
entire world is just not supportable at this time.''

\section{Obama Acts to Ease Flu Fears; U.S.~Says It Is Prepared}

WASHINGTON -- The Obama administration dispatched high-level officials from several agencies Monday
to allay concerns about swine flu and to demonstrate that it was fully prepared to confront the
outbreak even as the president said there was ``not a cause for alarm.''

Janet Napolitano, the secretary of homeland security, and Dr.~Richard E.~Besser, the acting director
of the Centers for Disease Control and Prevention, said the administration was prepared to respond
to any further spread of the swine flu virus.

Homeland security officials said they expected the outbreak to spread. ``We are proceeding as if we
are preparatory to a full pandemic level,'' Ms.~Napolitano said.

As the administration responds to its first domestic emergency, it is building on concrete
preparations made during the tenure of President George W.~Bush that have won praise from public
health experts. But its actions are also informed by what Mr.~Bush learned in his response to
Hurricane Katrina: that political management of a crisis, and of public expectations, can be as
important as the immediate response.

In a speech at the National Academy of Sciences on Monday, Mr.~Obama said only a few words about
swine flu. ``This is, obviously, a cause for concern and requires a heightened state of alert,'' he
said. ``But it's not a cause for alarm.''

But behind the scenes at the White House, aides said the president was directing his administration
to be ready in case an alarm needed to be sounded. A full report on the swine flu was added to Mr.
Obama's daily intelligence briefing, with updates given to him throughout the day.

Aides said they were mindful that how the president conducted himself in this period, both
substantively and stylistically, would be long remembered. But they adamantly rejected the idea that
this situation was at all comparable to that of the hurricanes that devastated the Gulf Coast in
2005.

Finding the right mix of alarm and reassurance is a delicate task for an elected official.

Eric Toner, a senior associate at the Center for Biosecurity at the University of Pittsburgh Medical
Center, said: ``It can be very dangerous to overreact. And it can be very dangerous to underreact.''
So far, Mr.~Toner said, Obama administration officials ``have managed to get it just right.''

Other public health experts also endorsed the administration's response to the outbreak that emerged
from Mexico. They gave much of the credit to President Bush, whose administration did extensive
planning for such an emergency.

``We're seeing a payoff of the original investment made in pandemic preparedness by the Bush
administration,'' said Jeffrey W.~Levi, executive director of Trust for America's Health. The term
pandemic refers to a widespread outbreak of an infectious disease.

Frances Fragos Townsend, who was assistant to President Bush for homeland security and
counterterrorism, noted that the Department of Health and Human Services had devised a detailed plan
for responding to the threat of pandemic flu in 2005 and 2006.

On his summer vacation in 2005, Mr.~Bush read ``The Great Influenza,'' a history of the 1918
pandemic by John M.~Barry.

Mr.~Obama also displayed interest in pandemic flu in 2005. Within months of taking office as a
senator, he introduced a bill to step up preparations, saying: ``We are in a race against time. The
nation's health officials have made some progress in preparing for pandemic influenza. Yet we have
much work to do.''

The swine flu outbreak has been linked to 149 deaths in Mexico. In the United States, the number of
confirmed cases grew to 50 on Monday.

While experts praised the Obama administration's initial response, many warned that a more extensive
outbreak of swine flu could tax the nation's public health capabilities.

``If this gets worse, you'll see the weakness of our system,'' said Dr.~Jeffrey P.~Koplan, a former
director of the Centers for Disease Control and Prevention. ``In an event like this, where
everyone's well-being is dependent on everyone else's, we will both feel and see the problems our
system creates.''

Dr.~Irwin Redlener, director of the National Center for Disaster Preparedness at Columbia
University, said federal officials reported six years ago that hospitals would need far more beds,
ventilators and personal protective equipment to respond to a pandemic. Hospitals never got nearly
enough extra equipment, Dr.~Redlener said.

``We will pay a very heavy price for this if we get the big one,'' he said.

The White House press secretary, Robert Gibbs, said Monday that the administration was not at a
disadvantage because of vacancies in top federal health positions.

``Our response is in no way hindered or hampered,'' Mr.~Gibbs said. When pressed to say whether
White House officials would prefer to have a full team in place, he said, ``We'd rather not have a
swine flu.''

Dr.~June E.~Osborn, former dean of the School of Public Health at the University of Michigan, said
the Obama administration appeared to be responding effectively, even without a secretary of health
and human services.

Mr.~Obama's nominee for health secretary, Gov. Kathleen Sebelius of Kansas, is waiting for Senate
confirmation. Senate action has been delayed in part by Republican questions about Ms.~Sebelius's
views on abortion.

On Tuesday, the Senate is scheduled to begin debate on the nomination, under an agreement that will
require 60 votes for confirmation.

Twenty positions at the Department of Health and Human Services are filled by the president, subject
to Senate confirmation. Mr.~Obama has nominated people to fill five of those positions, and none
have been confirmed.

Dr.~Besser, the acting director of the Centers for Disease Control, brings his experience as a past
director of the Office for Terrorism Preparedness and Emergency Response at the centers to his task
of dealing with the swine flu threat.

Ms.~Napolitano said her agency went into the crisis with more than a dozen vacancies in senior
positions, including the commissioner for customs and border patrol, the assistant secretary for
health affairs and the under secretary for intelligence analysis. She said those jobs were being
handled by career civil servants, working from detailed contingency plans inherited from the Bush
administration.

Dr.~Nicole Lurie, director of the public health preparedness program at the RAND Corporation, said,
``The federal government has come together with a pretty good, unified response'' to the swine flu
outbreak.

From the Hurricane Katrina experience, Dr.~Lurie said, federal officials learned ``the importance of
coordinating the government response, communicating with the public and mobilizing equipment'' as
fast as possible.

In recent years, she said, federal, state and local officials have conducted many exercises so they
would be prepared to respond to emergencies.

Since 2005, federal and state governments have spent more than \$1.5 billion to stockpile Tamiflu and
Relenza, antiviral medicines recommended by the government to treat infection with the swine flu
virus.

\section{Borders Are Tightened as Flu Cases Rise}

Countries around the world began tightening their border and immigration controls Tuesday as the
number of confirmed cases of swine flu continued to rise.

The number of deaths believed attributable to swine flu climbed to as many as 152 on Tuesday -- all
of them in Mexico -- as news agencies reported the number of confirmed cases of infection in the
United States stood at 50 after further testing at a New York City school.

Other cases have been reported in Ohio, Kansas, Texas and California. In addition, The Associated
Press reported that preliminary tests by health officials in New Jersey had identified five
``probable'' cases -- four people who were recently in Mexico and one who had been in California.

About 1,600 people in Mexico are now believed to have swine flu.

The World Health Organization raised its global pandemic flu alert level on Monday while
recommending that borders not be closed nor travel bans imposed. But in a possible precaution to be
taken by other nations, Japan said Tuesday it would no longer allow Mexican travelers to obtain a
visa upon arrival.

Health officers at three Japanese airports also were being deployed Tuesday to check passengers
before they disembarked from flights coming from Mexico, Canada and the United States. Travelers
suspected of having the flu would be quarantined and examined further at medical facilities, news
agencies reported, citing the Japanese health minister.

At least eight other countries in Asia were checking air passengers arriving from North America, and
China was tightening land border checks as well. Hong Kong, Taiwan and Russia were set to quarantine
passengers suspected of having the flu.

On Monday, Indonesia banned all pork products coming from Canada, the United States and Mexico, and
said it would destroy any pork products imported before the ban. International health experts said
the flu cannot be transmitted from pigs or from eating properly handled pork.

Russia and South Korea each reported a suspected case of swine flu on Tuesday. Tests were ongoing,
and South Korea raised its domestic alert level.

Two people in Scotland -- the first known victims of the virus in Britain -- were said by hospital
authorities on Tuesday to be recovering after contracting the flu while on honeymoon in Cancún,
Mexico. But the British authorities warned against nonessential travel and suggested that British
citizens in Mexico should consider leaving.

In an advisory on Tuesday, the British Foreign Office said: ``We are now advising against all but
essential travel to Mexico. Routine consular and all visa services at the embassy in Mexico City
have been suspended until further notice.''

It added: ``British nationals resident in or visiting Mexico may wish to consider whether they
should remain in Mexico at this time.''

The increase of the W.H.O. alert to Level 4 from Level 3 means that there has been sustained
human-to-human transmission. The change ``indicates that the likelihood of a pandemic has increased,
but not that a pandemic is inevitable,'' said Dr.~Margaret Chan, the W.H.O. director general.

The W.H.O.'s emergency committee, after meeting until 10:30 p.m. in Geneva, recommended abandoning
efforts to contain the flu's spread.

``Because the virus is already quite widespread in different locations, containment is not a
feasible option,'' said Dr.~Keiji Fukuda, the deputy director general.

The W.H.O. also recommended that vaccine makers keep making the seasonal flu vaccine instead of
switching over to a new one that matches the swine flu strain, but it urged them to start the
process of picking a pandemic strain, weakening it and making large batches of it, which could take
six months.

Dr.~Fukuda emphasized that the committee thought that ``a pandemic is not inevitable -- the
situation is fluid and will continue to evolve.''

In Mexico, state health authorities looking for the initial source of the outbreak toured a
million-pig hog farm in Perote, in Veracruz State. The plant is half-owned by Smithfield Foods, an
American company and the world's largest pork producer.

Mexico's first known swine flu case, which was later confirmed, was from Perote, according to Health
Minister Jos\'e \'Angel C\'ordova. The case involved a 5-year-old boy who recovered.

But a spokesman for the plant said the boy was not related to a plant worker, that none of its
workers were sick and that its hogs were vaccinated against flu.

American officials said their response to the epidemic was already aggressive, and the W.H.O.'s
decision to raise its pandemic alert to level 4 from level 3 would not change their plans. The
W.H.O. decision offered some official guidance to a world that, at least for the day, seemed swept
by confusion that unnerved international travelers and the financial markets. European and Asian
markets fell, and stock in airlines and the travel industry fell while those in pharmaceutical
companies rose.

Pharmacies in New York reported runs on Tamiflu, an antiflu drug -- something that public health
officials badly want to avoid because the drug could eventually be needed for the truly ill.

For now supplies of Tamiflu and Relenza, another antiflu drug, remain adequate, the manufacturers
said, but both were increasing production and expressed anxiety that shortages could develop if
governments placed huge orders.

The travel issue was the most confusing. On Monday, the European Union appeared to issue and then
rescind a ban on travel to the United States, drawing a rebuke from American officials, who
themselves later suggested that Americans drop all nonessential travel to Mexico.

None of the American cases have been serious, but Dr.~Richard Besser, acting director of the Centers
for Disease Control and Prevention, said he ``would not rest on that fact.''

``I expect that we will see additional cases, and I expect that the spectrum of disease will
expand,'' he said at a news conference.

Asked why the W.H.O. had waited so long to raise its alert level, Dr.~Fukuda said it was done on
technical grounds, that there was evidence of sustained human-to-human transmission of a new virus
and movement of that virus to new areas. But he conceded that ``the committee is very aware that
changes have quite significant political and economic effects on countries.''

The W.H.O. has no power to enforce any policies on member states, but different countries may have
their own pandemic flu plans that are triggered by changes in the alert level.

Spain and Scotland were the first countries outside North America to confirm cases. Suspected cases
appeared in Brazil, Australia, Israel and New Zealand, but confirmation is slow because most
nations' laboratories lack the test kit the C.D.C. is developing for the new virus.

The C.D.C. began sending out the new kits on Monday, meaning that soon some states and foreign
countries would be able to make their own diagnoses -- a development that could lead to a sharp
increase in confirmed cases.

Confusion regarding Europe's position on travel arose when the European Union's health commissioner,
Androulla Vassiliou, was questioned on a visit to Luxembourg and said Europeans ``should avoid
traveling to Mexico or the United States unless it's very urgent.''

Early reports of those remarks led both Dr.~Besser and Mayor Michael R.~Bloomberg of New York City
to publicly disagree.

``We don't think there is any reason not to travel and come to New York,'' the mayor said.

Ms.~Vassiliou's office later denied she had issued any travel advisory and said she was only
offering her personal opinion. ``She didn't want to insinuate risk where we're not sure,'' a
spokesman said, adding that formal advice would be offered later.

Mayor Bloomberg confirmed that there were now 28 cases in New York, all connected to St. Francis
Preparatory School in Fresh Meadows, Queens. He said there had been no suspected cases in any of the
city's intensive-care units. He acknowledged an increase in emergency-room visits, but he said his
preliminary information indicated that there were more people who were worried rather than seriously
ill.

New York's public health situation does not now resemble Mexico City's, the mayor said, and the
public does not need to don masks.

In Europe, a spokesman for the Stockholm-based European Center for Disease Prevention and Control
said there were about 40 suspected cases beyond the confirmed case in Spain.

The United States pork industry continued trying to allay consumer concerns about their products.
Many companies and hog farmers complained that the ``swine flu'' name was unfortunate and perhaps
inaccurate because, so far, the virus appears to be spreading without any contact with pigs.

``I guess everything has got to have a name,'' said Kyle Stephens, who raises show pigs in Amarillo,
Tex. ``The biggest thing we are up against is people thinking the worst, instead of checking into it
more.''

\section{Year After China Quake, New Births, Old Wounds}

MIANZHU, China -- Just 45 days old and swaddled in pink, Sang Ruifeng already has a purpose in life:
to bring to justice those responsible for the death of his 11-year-old brother.

Ruifeng will have to ensure, his father said, that the Chinese government gives a full accounting
for why thousands of students died in school collapses during the earthquake that devastated
southwest China one year ago. The brother that Ruifeng never knew was among 126 students crushed to
death in Fuxin No.~2 Primary School outside this lush farming town.

``I don't feel happy at all,'' the father, Sang Jun, said about the birth of his new son as his wife
bounced the baby up and down in a neighbor's home. ``I was telling my wife today, if we can't get
justice, we'll have our son carry on the quest for justice. This issue will be a burden on this
child.''

One year after the earthquake in Sichuan Province killed about 70,000 people and left 18,000
missing, mothers across the region are pregnant or giving birth again, aided by government medical
teams dispensing fertility advice and doing reverse-sterilization procedures. Because of China's
policy limiting most families to having one child, the students who died were often their parents'
only offspring. Officials say they hope a wave of births will help defuse the anger that many
grieving parents harbor over the collapses of so many schools on May 12, 2008, while nearby
buildings often remained standing.

But the wounds have festered, in part because the Chinese government, wary of any challenge to its
authoritarian rule, has muffled the parents and quashed public discussion of shoddy school
construction. As attention focuses again on Sichuan during the one-year anniversary of the
earthquake, the government has intensified its campaign to silence the parents and the media,
resorting to harassment by police and threats of imprisonment.

``The government says, 'Since you have a second child, why are you still asking about this?' '' said
Mr.~Sang, a former factory worker who was detained by the police in January when he tried to take a
train to Beijing to file a formal complaint. ``We tell the government: 'This is your responsibility,
this is your fault. So why shouldn't we question this?' ''

The Chinese government has refused to release the number of students who died or their names. But
one official report soon after the earthquake estimated that up to 10,000 students died in the
collapse of 7,000 classrooms and dormitory rooms.

Last year, officials in the central government announced that it would carry out an investigation
into the school collapses, but no results have been released. In March, a provincial official from
Sichuan told reporters in Beijing that the force of the earthquake rather than poor construction led
to the collapses.

On April 4, during Tomb-Sweeping Day, when Chinese honor the dead, groups of parents tried to gather
at the sites of collapsed schools to mourn their children. Plainclothes police officers quickly
surrounded them.

Propaganda officials recently ordered Chinese news organizations to report only positive
quake-related stories, and the Sichuan government has explicitly prohibited media organizations from
reporting on miscarriages by women in temporary housing camps. Some quake survivors say they fear
that the miscarriages may have been caused by high levels of formaldehyde in the prefabricated
housing.

``So many pregnant women are having miscarriages,'' said Ms.~Ren, a woman in a camp in Dujiangyan
who gave only her surname for fear of government reprisal. Her grandson was among the hundreds who
died at Xinjian Primary School. Her daughter-in-law got pregnant late last year, she said, but had a
miscarriage.

``The room in Dujiangyan People's Hospital was full of women with this same problem,'' Ms.~Ren said
as she wept.

The central government began sending fertility specialists to the earthquake region last year. The
Sichuan Daily, an official newspaper, reported on Feb.~29 that nearly 1,000 women in the quake
region had gotten pregnant, citing the Sichuan Province Family Planning Commission. Family planning
officials in Chengdu, the capital of Sichuan, declined requests for an interview.

One of the expectant mothers is Liu Li, whose brown corduroy overalls stretch across her swollen
belly. She said ``there were complicated feelings involved'' in finding out she was pregnant this
past winter.

``There was nervousness, happiness and also some guilt, because it was so soon after the death of
our first child,'' said Ms.~Liu, 35, as she filled a plastic basin in the Dujiangyan camp.

Like many parents whose children died, Ms.~Liu took a payment of about \$8,800 from the local
government and a guarantee of a pension in exchange for silence.

Many parents, even those expecting another child, have refused to quiet down.

Here in Mianzhu, women in more than half of the 126 households that lost children in Fuxin No.~2
Primary School are pregnant or gave birth recently, according to several parents. One father, Bi
Kaiwei, praised the free health care the government had provided for his wife, now more than four
months pregnant. But the pregnancy is no substitute for justice, he said.

Every day, the two visit the grave of their dead daughter. They have kept all her belongings,
including a stuffed white puppy and a light blanket now on the parents' bed. Framed photographs of
the girl are displayed throughout their prefabricated home just a few hundred yards from the site of
the Fuxin school.

``I feel this is the return of our daughter,'' Mr.~Bi's wife, Liu Xiaoying, said as she patted her
belly. ``But even though I'm comforting myself, telling myself this is her, I still don't feel
cheerful. I'm very depressed.''

Ms.~Liu was among a group of parents from Mianzhu who secretly traveled to Beijing in January to
file a petition with the central government. Officials there told them to file with the Sichuan
government.

But officials in Sichuan are trying to break the will of the parents. Mr.~Sang, the father with the
45-day-old son, said the police had threatened him with further imprisonment.

A man answering the telephone at the Mianzhu police station declined to comment.

On the edge of a wheat field here, Mr.~Sang has built a new home to replace the one that crumbled
during the earthquake. In one corner is a bedroom for his dead son, Xingpeng. Neatly stored inside
are a framed photograph of the boy and his most treasured possessions -- a fishing rod, white
dancing shoes, a glass fish tank.

The new son will not sleep here.

``We're going to keep this forever,'' Mr.~Sang said.

\section{Despite High Unemployment Numbers, Millions Still Being Hired}

Everyone knows the grim news -- unemployment in the United States has jumped to 8.5 percent, a
25-year high, and is racing toward double digits. Since November, the nation has lost more than
three million jobs.

But not everyone knows the brighter side to the equation: deep in the maw of the deepest recession
since the Great Depression, millions are still being hired.

So, while 4.8 million workers were laid off or chose to leave their jobs in February, employers
across the country hired 4.3 million workers that month, according to the Bureau of Labor
Statistics.

``The best thing you can say about these numbers is it speaks to the dynamism of the U.S.~economy,
and the net negative number that we all traffic in masks that,'' said Robert J.~Barbera, chief
economist at ITG, a research and trading firm. ``Ninety out of 100 people who know the number --
650,000 were lost in February -- think that means no one was hired and 650,000 were fired.''

In February -- before the economy started to show the first faint signs of a possible recovery --
there were three million job openings nationwide. And despite large new job losses likely to be
announced Friday, there are still millions of job openings.

Who is hiring? Hospitals, colleges, discount stores, restaurants and municipal public works
departments. I.B.M. is hiring more than 700 people for its new technical services center in Dubuque,
Iowa, while the Cleveland Clinic has 500 job openings, not just for nurses but also for pharmacy
aides and physical therapists. And after President Obama's stimulus package kicks into gear, state,
local governments and road-building contractors are expected to increase hiring.

Zachary Schaefer has hired 72 people since February for the Culver's hamburger and frozen custard
restaurant that he and several partners just opened in Surprise, Ariz.

``The amount of applicants who are qualified is definitely up,'' compared with previous hiring
efforts, he said. ``Whereas before we were counting on a lot of high school applicants, now there
are a lot more middle-age people applying.''

Eddie Hamm, a former construction worker, was unemployed for five months when he drove by the site
where the Culver's was under construction. Mr.~Hamm, 29, applied for a job there, and now he's a
``fry guy.''

``I'm just happy I got hired -- I didn't want to stay home, not doing anything,'' he said, hardly
complaining that he is earning half the \$15 an hour he made in construction. ``I don't look at it
like I'm making \$7.50. I look at it -- I'm having a job in a down time, and it's a job where I can
move up.''

Economists and job counselors advise the unemployed that there are definitely jobs to be had, even
if there aren't nearly enough to go around. With 13.2 million people out of work, there are 4 1/3
unemployed Americans for every job opening. ``You're facing more competition for every job you apply
for, but the reality is there is a lot of hiring going on,'' said Andrew M.~Sum, director of the
Center for Labor Market Studies at Northeastern University. ``You're never going to find anything
unless you apply.''

Even industries that have taken a beating are doing plenty of hiring. According to the Bureau of
Labor Statistics, construction companies hired 366,000 workers in February, and manufacturers hired
249,000. Retailers hired 536,000 workers in February, but that was down 25 percent from the previous
February.

Some job openings are to replace retirees, some to replace employees who left for other jobs, but
many openings result from expansion. Companies that are still growing are blessed with talented
applicants.

``It's easier to hire in a recession -- we have about five applications for every position,'' said
Howard Glickberg, principal owner of Fairway Market, the well-known grocery company based in
Manhattan.

Fairway just hired 350 people for its month-old store in Paramus, N.J., the first Fairway outside of
New York State. The company plans to add 1,200 more workers over the next two years by opening
stores in Queens; Pelham Manor, N.Y.; and Stamford, Conn.

``What you have to be afraid of is hiring someone who can't find something better at the time, and
when they find something better they leave you,'' Mr.~Glickberg said. ``I want to hire someone who
will make a career of it.''

The nation's largest private-sector employer, Wal-Mart Stores, is also hiring aplenty. Wal-Mart,
with 1.4 million workers nationwide, hires several hundred thousand workers each year because of
employee turnover, and expects to increase its domestic work force by nearly 50,000 this year,
thanks to plans to open 150 new stores.

Shawnalyn Conner is running a hiring center for a Wal-Mart store that will open on June 17 in
Weaverville, N.C., near Asheville. She plans to hire 350 workers.

``The biggest comment that we get from people is that they're looking for a company that's growing,
and Wal-Mart offers that,'' said Ms.~Conner, who, as the top manager of the new store, has hired 77
people so far. Gisel Ruiz, senior vice president for the people division of Wal-Mart U.S., said the
company had a hiring program for former junior military officers, often for jobs as assistant store
managers. With many veterans having a hard time landing jobs, Wal-Mart hired 150 former officers
last year.

The health care industry has held its own in hiring. The University of Miami medical school, which
runs three hospitals, has 250 openings and is hiring about 35 people a month, compared with 100 a
month in good times. Cleveland Clinic has 500 job openings, compared with 2,000 during better times.

``We have a hiring freeze on, but even when there's a hiring freeze, we need to maintain our head
count,'' said Joe Patrnchak , Cleveland Clinic's chief human resources officer. ``We have 40,000
people, and you're going to have some openings.''

He is encountering an unusual snag in hiring people. ``A challenge we have now is people from other
areas are having problems selling their homes,'' Mr.~Patrnchak said. ``People aren't quite as mobile
nowadays.''

The University of Miami medical school is also facing an unexpected problem. ``There's a flood of
applicants, but even so, it's harder to find really good, experienced people,'' said Paul Hudgins,
its associate vice president for medical human resources. ``We're seeing people hunkering down and
saying they're going to stay where they are.''

The recession has encouraged people to cling to their jobs. Just 1.5 percent of workers voluntarily
quit their jobs in February, the lowest level since the Bureau of Labor Statistics began collecting
those numbers eight years ago.

Like many educational institutions, Washington University in St. Louis continues to hire. It has 175
job openings in admissions, residential life and other areas. There is a flood of job applicants,
and Ann Prenatt, vice chancellor for human resources, said that has pros and cons, the advantage
being that the university does not have to offer large premiums as often to draw coveted applicants.

As for the disadvantage, Ms.~Prenatt said: ``People are applying for many positions in the hope that
something will stick even though they don't remotely meet the minimum qualifications. That can be
very frustrating for the applicant, but it also creates a logjam for us.''

\section{Obama's Plan on Corporate Taxes Unnerves the Indian Outsourcing Industry}

NEW DELHI -- President Obama's proposal to change the American corporate tax system is winning few
fans in India, where some say it is aimed at curbing the country's outsourcing industry.

Perhaps that is because Mr.~Obama specifically struck out at the epicenter of Indian outsourcing.

The president vowed Monday to overhaul a tax code that allowed companies to pay less tax, as he put
it, to ``create a job in Bangalore, India, than if you create one in Buffalo, New York.'' One
element to that change could be the elimination of a deduction for American companies when they
invest in subsidiaries outside the United States.

American companies have tens of thousands of employees in India in wholly owned subsidiaries. Many
of these Indian operations handle customer service and back-office functions, particularly for banks
and credit card companies. American businesses employ thousands more in India by contracting work to
local technology and outsourcing companies.

And recently, many American corporations have also expanded their sales, marketing and distribution
in India to take advantage of the country's fast economic growth and expanding middle class.

Many business people in India were upset by Mr.~Obama's tax proposal. The president of the
Associated Chambers of Commerce and Industry of India, Sajjan Jindal, said it could ``kill the
spirit of competition.''

The Indian affiliate of CNN spent Tuesday afternoon asking economists and politicians whether
Mr.~Obama was ``anti-India.'' An editorial in The Times of India said Bangalore had become a
``catch-all term to hang U.S.~economic woes on.''

What is unclear, though, is what, if any, impact Mr.~Obama's proposed tax plan will actually have on
jobs in India.

``It's a tax disincentive to discourage outsourcing to countries like India,'' said Uday Ved, head
of tax issues at KPMG India. But according to Mr.~Ved and other international tax experts, companies
do not move jobs to India because the tax rate is lower; they do it because labor costs less.

``We still believe that the cost advantage to India is so high'' that American companies will
continue to move some jobs to India, Mr.~Ved said.

Raymond J.~Wiacek, chairman of the global tax practice at Jones Day, said, ``I don't think it's
going to make a bean's difference to India.''

He added, ``India is a highly skilled but inexpensive labor market,'' and not one where American
companies have been accruing enormous profits in their foreign subsidiaries.

Some big American companies have large numbers of employees in India. For example, General Electric
has about 14,500 employees in India, I.B.M. more than 74,000, and Citigroup more than 10,000. In
addition, India's information technology and outsourcing companies employ about 2.2 million people,
and American companies account for about 60 percent of their business.

``The jobs aren't coming back to the U.S.~as a result of this proposal,'' Mr.~Wiacek said. The tax
proposal is ``about revenues and that's it.''

\pagebreak
\printindex
\end{document}
